\documentclass[11pt,oneside,openany]{book}

\usepackage[margin=1in,headheight=15pt]{geometry}
\usepackage[T1]{fontenc}
\usepackage{lmodern}
\usepackage{microtype}
\usepackage{mathtools,amssymb,amsthm}
\usepackage{tikz}
\usetikzlibrary{arrows.meta,calc,patterns,positioning,shapes.geometric}
\usepackage[most]{tcolorbox}
\usepackage{enumitem}
\usepackage{needspace}
\usepackage{booktabs,tabularx,array,multicol}
\usepackage{xcolor}
\usepackage{fancyhdr}
\usepackage{hyperref}

\definecolor{ink}{HTML}{17243A}
\definecolor{indigo}{HTML}{293D70}
\definecolor{violet}{HTML}{6750A4}
\definecolor{mint}{HTML}{2A8C82}
\definecolor{amber}{HTML}{E9A23B}
\definecolor{rose}{HTML}{D85D75}
\definecolor{mist}{HTML}{EAF5F2}
\definecolor{paper}{HTML}{FCFAF5}
\definecolor{muted}{HTML}{647084}

\pagecolor{paper}
\color{ink}
\hypersetup{colorlinks=true,linkcolor=violet,urlcolor=rose,citecolor=mint,pdfauthor={Course design by OpenAI Codex},pdftitle={Chance, Data, and Decisions}}
\setlength{\parindent}{0pt}
\setlength{\parskip}{5pt}
\setlist{itemsep=2pt,topsep=5pt}
\setlist[enumerate,1]{label=\textbf{\arabic*.},leftmargin=*}

\pagestyle{fancy}
\fancyhf{}
\fancyhead[L]{\small\color{muted}\nouppercase{\leftmark}}
\fancyhead[R]{\small\color{muted}Chance, Data, and Decisions}
\fancyfoot[C]{\color{violet}\thepage}
\renewcommand{\headrulewidth}{0.3pt}
\renewcommand{\chaptermark}[1]{\markboth{#1}{}}

\newtheoremstyle{course}{6pt}{6pt}{\itshape}{}{\bfseries\color{indigo}}{.}{.5em}{}
\theoremstyle{course}
\newtheorem{theorem}{Theorem}[chapter]
\newtheorem{lemma}[theorem]{Lemma}
\newtheorem{proposition}[theorem]{Proposition}
\newtheorem{corollary}[theorem]{Corollary}
\theoremstyle{definition}
\newtheorem{definition}[theorem]{Definition}
\newtheorem{example}[theorem]{Example}
\newtheorem{question}[theorem]{Guiding question}

\newtcolorbox{bigidea}[1][]{enhanced,breakable,colback=mist,colframe=mint,boxrule=.8pt,arc=2mm,left=5mm,borderline west={3pt}{0pt}{mint},title={Big idea},fonttitle=\bfseries,#1}
\newtcolorbox{checkpoint}[1][]{enhanced,breakable,colback=white,colframe=amber,boxrule=.7pt,arc=2mm,title={Checkpoint},fonttitle=\bfseries\color{ink},#1}
\newtcolorbox{warningbox}[1][]{enhanced,breakable,colback=rose!5,colframe=rose,boxrule=.7pt,arc=2mm,title={Common trap},fonttitle=\bfseries,#1}
\newtcolorbox{labbox}[1][]{enhanced,breakable,colback=indigo!4,colframe=indigo,boxrule=.7pt,arc=2mm,title={Explore},fonttitle=\bfseries,#1}
\newtcolorbox{exercisebox}[1][]{enhanced,breakable,colback=white,colframe=indigo!45,boxrule=.6pt,arc=1mm,title={Exercises},fonttitle=\bfseries\color{indigo},#1}

\newcommand{\R}{\mathbb R}
\newcommand{\N}{\mathbb N}
\newcommand{\E}{\mathbb E}
\renewcommand{\P}{\mathbb P}
\newcommand{\Var}{\operatorname{Var}}
\newcommand{\Cov}{\operatorname{Cov}}
\newcommand{\Bin}{\operatorname{Bin}}
\newcommand{\Pois}{\operatorname{Pois}}
\newcommand{\Normal}{\mathcal N}
\newcommand{\given}{\,\middle|\,}
\newcommand{\ind}{\mathbf 1}
\newcommand{\exlevel}[1]{\hfill\textcolor{violet}{\small[#1]}}
\newcommand{\chaptertag}[1]{\par\vspace{-1.4em}{\color{violet}\large\sffamily #1}\par\vspace{1em}}

\begin{document}

\hypersetup{pageanchor=false}
\begin{titlepage}
\begin{tikzpicture}[remember picture,overlay]
  \fill[indigo] (current page.south west) rectangle (current page.north east);
  \draw[mint!70,line width=1.4pt] plot[smooth] coordinates {
    ($(current page.south west)+(1.0,4.0)$)
    ($(current page.south west)+(2.4,4.2)$)
    ($(current page.south west)+(3.7,5.3)$)
    ($(current page.south west)+(5.2,8.2)$)
    ($(current page.south west)+(6.7,5.3)$)
    ($(current page.south west)+(8.0,4.2)$)
    ($(current page.south west)+(9.4,4.0)$)};
  \foreach \x/\y in {2.0/10.2,2.6/11.1,3.2/9.8,3.8/12.0,4.4/10.7,5.0/11.8,5.6/10.0,6.2/12.4,6.8/11.2,7.4/10.5}{
    \fill[amber] ($(current page.south west)+(\x,\y)$) circle (2.3pt);
  }
  \draw[rose,line width=3pt] ($(current page.south west)+(0,2.0)$) -- ($(current page.south east)+(0,2.0)$);
\end{tikzpicture}
\vspace*{4.1cm}
{\color{amber}\sffamily\bfseries\Large AN INTRODUCTION TO PROBABILITY AND STATISTICS\par}
\vspace{.5cm}
{\color{white}\fontsize{36}{42}\selectfont\bfseries Chance, Data,\par and Decisions\par}
\vspace{1cm}
{\color{white!78}\Large A visual, inquiry-driven course with exercises\par}
\vfill
{\color{white!75}\large One-semester course packet\par}
{\color{white!55}\normalsize events $\to$ random variables $\to$ samples $\to$ inference\par}
\end{titlepage}

\hypersetup{pageanchor=true}
\frontmatter
\chapter*{How to use this course}
\addcontentsline{toc}{chapter}{How to use this course}

\begin{bigidea}
Probability starts with a model of what might happen; statistics works backward from what did happen. The course repeatedly follows one cycle:
\[
\text{model} \longrightarrow \text{simulate} \longrightarrow \text{observe} \longrightarrow \text{infer} \longrightarrow \text{check}.
\]
The point is not merely to calculate a probability or a confidence interval, but to understand which assumptions made that calculation meaningful.
\end{bigidea}

\section*{Audience and prerequisites}

This packet is designed for a first undergraduate course or a mathematically ambitious secondary course. Students should be comfortable with algebra, functions, summation notation, and basic differentiation and integration. Multivariable calculus is not required. Proofs are included, but the emphasis is on modeling, interpretation, and the disciplined use of computation.

\section*{Learning goals}

By the end, students should be able to:
\begin{enumerate}
  \item translate verbal questions into sets, events, and probability models;
  \item use counting, conditioning, independence, and Bayes' rule correctly;
  \item distinguish permutations from combinations, prove basic binomial identities, and apply inclusion--exclusion to both counts and probabilities;
  \item derive the birthday-paradox formula from an explicit probability space and distinguish any-pair matches from matches with one designated person;
  \item move fluently among CDFs, probability masses, and densities;
  \item compute and interpret expectation, variance, covariance, and correlation;
  \item recognize standard distributions and justify an appropriate model;
  \item obtain marginal and conditional distributions from a joint model;
  \item explain the law of large numbers and central limit theorem through both mathematics and simulation;
  \item evaluate estimators by bias, standard error, and mean squared error;
  \item construct and interpret basic confidence intervals and hypothesis tests;
  \item perform a simple Bayesian update and compare credible and confidence intervals;
  \item distinguish exchangeability from independence and explain how a shared random parameter produces dependence and a random limiting frequency.
\end{enumerate}

\section*{A 14-week route}

The design assumes two 75-minute meetings per week. ``Lab'' means that students make a prediction, run an experiment or simulation, and explain a discrepancy.

\begin{center}
\small
\begin{tabularx}{\textwidth}{>{\bfseries}p{.07\textwidth}p{.41\textwidth}X}
\toprule
Week & Core mathematics & Suggested lab or assessment \\
\midrule
1 & Sets, events, Venn diagrams, counting & Inclusion--exclusion gallery \\
2 & Probability spaces and axioms & Build and critique sample spaces \\
3 & Conditional probability and trees & Medical-test simulation \\
4 & Independence, total probability, Bayes & Monty Hall and base-rate lab \\
5 & Random variables, CDFs, PMFs, densities & One distribution in three representations \\
6 & Expectation, variance, moments & Indicator-variable studio \\
7 & Bernoulli, binomial, geometric, Poisson & Rare-event approximation lab \\
8 & Uniform, exponential, normal, beta & Memorylessness and normal areas \\
9 & Joint, marginal, conditional laws & Heat maps and independence \\
10 & Covariance, correlation, conditional expectation & Correlation counterexamples \\
11 & Laws of large numbers; Monte Carlo & Running-average experiment \\
12 & Central limit theorem; sampling distributions & Visual CLT; Problem Set 5 \\
13 & Estimation, confidence intervals, tests & Coverage and power simulation \\
14 & Bayesian inference and synthesis & Prior-to-posterior capstone \\
\bottomrule
\end{tabularx}
\end{center}

\section*{Assessment design}

\textbf{Optional synthesis.} Chapter~\ref{ch:exchangeability} follows Bayesian inference as an exchangeability capstone. Its coin and urn examples can anchor a week-14 synthesis session; the optional proof is extension reading, not an added prerequisite for the first weeks. A short preview in Chapter~3 plants the distinction between symmetry and independence.

\textbf{Early-course companion.} Chapters 1 and 2 now provide a fuller
development of permutations, combinations, binomial coefficients, both forms
of inclusion--exclusion, finite probability spaces, and the birthday paradox.
The companion Problem Set 2 revisits the probability spaces introduced in
week two and connects them to this counting material. Its worked solutions
are supplied separately for the instructor.

One workable balance is 35\% problem sets, 20\% short quizzes, 20\% computational labs, and 25\% a modeling capstone. Exercises use three labels:
\begin{itemize}
  \item \textbf{Fluency}: calculations and translations that make later reasoning automatic.
  \item \textbf{Reasoning}: derivations, short proofs, and careful interpretation.
  \item \textbf{Exploration}: simulation, visualization, model criticism, or open-ended comparison.
\end{itemize}

\section*{Dependency map}

\begin{center}
\begin{tikzpicture}[
  node distance=8mm and 10mm,
  box/.style={draw=indigo!55,rounded corners=2mm,fill=white,minimum height=9mm,align=center,font=\small,inner xsep=4mm},
  arr/.style={-{Stealth[length=2mm]},thick,mint}
]
\node[box] (sets) {Sets and\\counting};
\node[box,right=of sets] (spaces) {Probability\\spaces};
\node[box,right=of spaces] (cond) {Conditioning\\and Bayes};
\node[box,below=of spaces] (rv) {Random\\variables};
\node[box,right=of rv] (mom) {Moments and\\standard laws};
\node[box,below=of rv] (joint) {Joint\\distributions};
\node[box,right=of joint] (limit) {LLN and\\CLT};
\node[box,right=of limit] (infer) {Sampling and\\inference};
\node[box,below=of infer] (bayes) {Bayesian\\inference};
\draw[arr] (sets)--(spaces);
\draw[arr] (spaces)--(cond);
\draw[arr] (spaces)--(rv);
\draw[arr] (rv)--(mom);
\draw[arr] (rv)--(joint);
\draw[arr] (mom)--(joint);
\draw[arr] (joint)--(limit);
\draw[arr] (limit)--(infer);
\draw[arr] (cond) to[bend left=18] (bayes);
\draw[arr] (infer)--(bayes);
\end{tikzpicture}
\end{center}

\begin{checkpoint}
\textbf{A good opening prompt.} A family has two children and at least one is a girl. What is the probability that both are girls? Ask students to defend not just an answer, but a sample space. Then change how the information was obtained and ask whether the answer changes.
\end{checkpoint}

{\setlength{\parskip}{0pt}\tableofcontents}

\mainmatter
\chapter{The grammar of uncertainty}
\label{ch:counting}
\chaptertag{Sets, permutations, combinations, and inclusion--exclusion}

\begin{question}
Before assigning probabilities, how do we describe precisely what could happen?
\end{question}

A \emph{sample space} is built from possible outcomes; an \emph{event} is a collection of outcomes. Set notation is therefore the grammar of probability.

\section{Set operations}

Let $\Omega$ be a universe of outcomes and $A,B\subseteq\Omega$. The union $A\cup B$ means ``$A$ or $B$ (or both),'' the intersection $A\cap B$ means ``both,'' and the complement $A^c$ means ``not $A$.'' The difference $A\setminus B$ contains outcomes in $A$ but not $B$.

\begin{center}
\begin{tikzpicture}[scale=.95]
  \draw[rounded corners,thick,indigo] (-3,-1.6) rectangle (3,1.8);
  \fill[mint!28] (-.9,.1) circle (1.25);
  \fill[amber!32] (.9,.1) circle (1.25);
  \draw[mint!90!black,thick] (-.9,.1) circle (1.25);
  \draw[amber!85!black,thick] (.9,.1) circle (1.25);
  \node at (-1.25,.25) {$A$};
  \node at (1.25,.25) {$B$};
  \node[above right] at (-3,1.8) {$\Omega$};
  \node[below] at (0,-1.8) {Overlap represents $A\cap B$};
\end{tikzpicture}
\end{center}

De Morgan's laws translate negated statements:
\[
(A\cup B)^c=A^c\cap B^c,
\qquad
(A\cap B)^c=A^c\cup B^c.
\]

\begin{warningbox}
In everyday speech, ``or'' is sometimes exclusive. In probability, $A\cup B$ is inclusive unless explicitly stated otherwise.
\end{warningbox}

\section{Products and sequences}

The Cartesian product $A\times B$ contains ordered pairs $(a,b)$. Repeated trials naturally live in product spaces. Three coin flips have sample space $\{H,T\}^3$, with eight ordered strings. Order matters: $HTT$ and $THT$ are different outcomes.

\section{The multiplication principle}

Start by deciding what one outcome records. An ordered list, an unordered
group, and a list with repetitions are different objects. Before using a
formula, describe when two proposed outcomes count as the same.

The \emph{addition principle} says that counts add over disjoint cases.
For example, a club offers either one of three morning sessions or one of
five afternoon sessions. Choosing exactly one session gives $3+5=8$
possibilities. Choosing one session in \emph{each} period gives $3\cdot5=15$.
The words ``either'' and ``each'' describe different experiments.

If a first choice has $m$ possibilities and, after \emph{each} first choice,
a second choice has $n$ possibilities, there are $mn$ ordered outcomes.
More generally, if the successive numbers of choices are $n_1,\ldots,n_k$,
the number of outcomes is $n_1\cdots n_k$. The available choices may depend
on earlier choices; their \emph{number} at each stage must be the stated one.
If different first choices leave different numbers of continuations, count
each branch and add its contributions instead.

\begin{example}[Codes and officers]
A four-digit code may begin with zero. If digits can repeat, there are
$10^4$ codes. If digits must be distinct, the count is
$10\cdot9\cdot8\cdot7=5040$. Assigning president, secretary, and treasurer
from ten students similarly gives $10\cdot9\cdot8=720$ outcomes when no
student holds two offices.
\end{example}

\section{Permutations and combinations}

There are $n!=n(n-1)\cdots1$ orders of $n$ distinct objects, with $0!=1$.
For $0\le k\le n$, an ordered selection of $k$ distinct objects from $n$
has
\[
  (n)_k=n(n-1)\cdots(n-k+1)=\frac{n!}{(n-k)!}
\]
possibilities. This is sometimes written $P(n,k)$. An unordered selection
of $k$ objects is a \emph{combination}. Every such selection has exactly
$k!$ orderings, so
\[
  \binom nk=\frac{(n)_k}{k!}=\frac{n!}{k!(n-k)!}.
\]
In particular, $\binom n0=\binom nn=1$: there is one empty selection and
one selection containing everything. We use $\binom nk=0$ when $k<0$ or
$k>n$ and $n$ is a nonnegative integer.

\begin{center}
\small
\begin{tabularx}{\textwidth}{@{}XXl@{}}
\toprule
Experiment & What distinguishes outcomes? & Count \\
\midrule
$k$ draws with replacement & Order matters; repetition allowed & $n^k$ \\
$k$ draws without replacement & Order matters; objects distinct & $(n)_k$ \\
A $k$-element subset & Order ignored; no repetition & $\binom nk$ \\
An unordered selection with repetition & Only the number of each type matters & $\binom{n+k-1}{k}$ \\
\bottomrule
\end{tabularx}
\end{center}

\begin{example}[A committee is not a list of officers]
From ten students there are $\binom{10}{3}=120$ three-person committees,
but $720$ assignments to three distinct offices. A committee with one
designated chair has $3\binom{10}{3}=360$ possibilities: the two other
members have no distinct roles. The question about order must be answered
before choosing a formula.
\end{example}

\begin{proposition}[Arrangements with repeated objects]
If $n$ objects consist of $r$ types with multiplicities
$n_1,\ldots,n_r$, where $n_1+\cdots+n_r=n$, the number of distinct
arrangements is $n!/(n_1!\cdots n_r!)$.
\end{proposition}
\begin{proof}
Temporarily distinguish objects of the same type with labels. There are
$n!$ arrangements of the labeled objects. Each arrangement after erasing
the labels comes from exactly $n_1!\cdots n_r!$ labeled arrangements.
Dividing by that constant multiplicity gives the result.
\end{proof}
For example, $\texttt{BANANA}$ has $6!/(3!2!)=60$ distinct arrangements.

\subsection*{Unordered selections with repetition}
Suppose $k$ scoops are chosen from $n\ge1$ flavors, with repetition allowed
and order ignored. An outcome specifies nonnegative counts
$x_1+\cdots+x_n=k$. Arrange $k$ stars and $n-1$ bars in a row: stars
before the first bar count flavor 1, those between the first two bars
count flavor 2, and so on. Empty groups are allowed. Choosing the $k$
star positions among $n+k-1$ positions gives $\binom{n+k-1}{k}$ outcomes.
For two scoops from three flavors $A,B,C$, these are
$AA,AB,AC,BB,BC,CC$, a total of $\binom42=6$.

\begin{warningbox}
Divide a count only when every desired outcome has the same number of
representations. Ordered pairs $AB,BA$ both become $AB$, but $AA$ has
only one ordering. Thus dividing $3^2$ by $2!$ does not count the six
two-scoop selections. The equal-multiplicity argument for combinations
applies to distinct objects.
\end{warningbox}

\section{Binomial coefficients count in two directions}

\begin{proposition}[Symmetry and Pascal's identity]
For $0\le k\le n$,
\[
\binom nk=\binom n{n-k},
\]
and, for $n\ge1$,
\[
\binom nk=\binom{n-1}{k-1}+\binom{n-1}{k}.
\]
\end{proposition}
\begin{proof}
A $k$-element subset determines its $(n-k)$-element complement, proving
symmetry. For Pascal's identity, single out one object. A $k$-element
subset either contains it, leaving $k-1$ choices from the other $n-1$
objects, or omits it, leaving $k$ choices from those objects. These two
cases are disjoint and exhaustive.
\end{proof}

\begin{theorem}[Binomial theorem]
For a nonnegative integer $n$,
\[
  (x+y)^n=\sum_{k=0}^n\binom nk x^{n-k}y^k.
\]
\end{theorem}
\begin{proof}
Expand the product of $n$ copies of $(x+y)$. To produce $x^{n-k}y^k$,
choose the $k$ factors that supply $y$; all the others supply $x$.
There are $\binom nk$ choices.
\end{proof}

\begin{example}[Coefficients carry both counts and weights]
In $(2-x)^4$, choosing $k$ copies of $-x$ and $4-k$ copies of $2$ gives
\[
 (2-x)^4=16-32x+24x^2-8x^3+x^4.
\]
For instance, the $x^2$ coefficient is $\binom42 2^2(-1)^2=24$.
The binomial coefficient counts the positions; the remaining factors
give the contribution from each choice.
\end{example}

The same distinction explains a probability formula. In $n$ independent
trials with success probability $p$, a \emph{specified} string with $k$
successes has probability $p^k(1-p)^{n-k}$. There are $\binom nk$ such
strings, so the probability of exactly $k$ successes is
$\binom nk p^k(1-p)^{n-k}$. Their sum is
$(p+(1-p))^n=1$. Chapter 5 develops this binomial distribution; for
$p=0$ or $1$, every trial has its certain outcome.

Taking $x=y=1$ gives $\sum_{k=0}^n\binom nk=2^n$. This also counts all
subsets in two ways: by their sizes, or by making an include/exclude
choice for each object. Taking $x=1,y=-1$ gives, for $n\ge1$,
$\sum_{k=0}^n(-1)^k\binom nk=0$. This alternating sum is the engine
behind the general inclusion--exclusion formula.

\section{Inclusion--exclusion for finite sets}

\begin{proposition}[Two and three sets]
For finite sets $A,B,C$,
\begin{align*}
|A\cup B|&=|A|+|B|-|A\cap B|,\\
|A\cup B\cup C|&=|A|+|B|+|C|\\
&\quad-|A\cap B|-|A\cap C|-|B\cap C|
  +|A\cap B\cap C|.
\end{align*}
\end{proposition}

For two sets, subtraction repairs the double count of the overlap. With
three sets, an element in all three is added three times and subtracted
three times, so it must be added back once. An element in exactly two
sets contributes $2-1=1$; an element in exactly one contributes $1$.
An element outside the union contributes zero. This proves the formulas.

\begin{warningbox}
$A\cap B$ includes the triple intersection $A\cap B\cap C$.
It does not mean ``in $A$ and $B$ but outside $C$.'' Confusing these
two regions produces incorrect inclusion--exclusion calculations.
\end{warningbox}

\begin{example}[Three activities]
Among 100 students, 39 attend activity $A$, 29 attend $B$, and 25 attend
$C$. The pairwise intersection counts are 13, 11, and 9, respectively,
and 5 students attend all three. Thus the union has
$39+29+25-13-11-9+5=65$ students; 35 attend none.
The exactly-two count is $(13-5)+(11-5)+(9-5)=18$.
The exactly-one count is $65-18-5=42$.
\end{example}

\begin{theorem}[Finite inclusion--exclusion]
For finite sets $A_1,\ldots,A_r$,
\[
\left|\bigcup_{i=1}^r A_i\right|
=\sum_{\varnothing\ne J\subseteq\{1,\ldots,r\}}
 (-1)^{|J|+1}\left|\bigcap_{j\in J}A_j\right|.
\]
Equivalently: add the single-set sizes, subtract the pairwise
intersection sizes, add the triple-intersection sizes, and continue
with alternating signs through the intersection of all $r$ sets.
\end{theorem}
\begin{proof}
An element belonging to exactly $t\ge1$ sets occurs in
$\binom tk$ of the $k$-fold intersections. Its total contribution is
\[
\sum_{k=1}^t(-1)^{k+1}\binom tk=1-(1-1)^t=1.
\]
An element in no set contributes zero. Thus each element of the union
is counted exactly once.
\end{proof}

\begin{example}[Using all three symbols]
There are $3^4$ strings of length four over $\{A,B,C\}$.
To count those using every symbol, subtract strings missing at least
one symbol. Each missing-symbol event has $2^4$ strings; each pair of
missing-symbol events has one string; no length-four string misses
all three symbols. Therefore the required count is
\[
3^4-3\cdot2^4+3\cdot1^4=36.
\]
These are counts; turning them into probabilities requires a specified
probability model.
\end{example}

\begin{example}
A five-card poker hand is an unordered choice from 52 cards, so there are $\binom{52}{5}$ hands. Exactly two aces means choosing two of four aces and three of the remaining 48 cards:
\[
\binom42\binom{48}{3}.
\]
\end{example}

\begin{labbox}
Use the website's counting lab to compare committees, ordered selections,
and repeated draws. Then use the inclusion--exclusion lab to vary the eight
disjoint regions of three sets, including the region outside the union.
Recover the marginal and intersection totals from those regions. With
all seven marginal/intersection totals specified, the seven interior
region counts are determined; knowing the universe size also determines
the outside region. If totals are omitted, the remaining data need not
determine a unique diagram.
\end{labbox}

\begin{exercisebox}
\begin{enumerate}
  \item Write $(A\cap B^c)\cup(A^c\cap B)$ in words. \exlevel{Fluency}
  \item Verify both De Morgan laws using membership tables. \exlevel{Reasoning}
  \item How many eight-character strings contain exactly three letters and five digits? Use 26 letters and 10 digits, allow repetition, and allow a leading zero. \exlevel{Fluency}
  \item Count five-card hands containing at least one ace. \exlevel{Reasoning}
  \item Prove $\binom nk=\binom n{n-k}$ by a counting argument. \exlevel{Reasoning}
  \item Among integers from 1 to 1,000, how many are divisible by 2 or 3 or 5? \exlevel{Reasoning}
  \item Design two different sampling procedures that produce the same set of possible outcomes but not the same probabilities. \exlevel{Exploration}
  \item From nine students, count a four-person committee, a list of four distinct officers, and a four-person committee with one chair. \exlevel{Fluency}
  \item Count the distinct arrangements of the letters in $\texttt{LEVEL}$. Explain the divisor in your formula. \exlevel{Fluency}
  \item Prove Pascal's identity by deciding whether a distinguished object is selected. \exlevel{Reasoning}
  \item Find the coefficient of $x^2$ in $(1+3x)^5$, and explain its binomial coefficient. \exlevel{Fluency}
  \item Count length-five strings over $\{A,B,C\}$ that use all three symbols. \exlevel{Reasoning}
  \item In five independent trials with success probability $1/3$, find
  the probability of exactly three successes. Distinguish one specified
  string from all strings with three successes. \exlevel{Fluency}
  \item Choose four scoops from three flavors, allowing repetition.
  Count ordered lists and unordered selections. Explain why dividing
  the first count by $4!$ is invalid. \exlevel{Reasoning}
  \item Count the distinct arrangements of $\texttt{MISSISSIPPI}$.
  Identify every repeated-letter factor. \exlevel{Fluency}
  \item Prove $\sum_{k=0}^n k\binom nk=n2^{n-1}$ for $n\ge1$ by
  counting subsets with one distinguished member in two ways.
  Appendix~\ref{app:calculus} also obtains this identity by differentiation.
  \exlevel{Reasoning}
\end{enumerate}
\end{exercisebox}

\chapter{A measure of possibility}
\chaptertag{Probability spaces, inclusion--exclusion, and the birthday paradox}

\begin{question}
What rules must any coherent numerical description of uncertainty obey?
\end{question}

\section{The axioms}

\begin{definition}
A \emph{probability space} is a triple $(\Omega,\mathcal F,\P)$: a sample space $\Omega$, a collection $\mathcal F$ of events containing $\Omega$ and closed under complements and countable unions, and a function $\P:\mathcal F\to[0,1]$ satisfying
\begin{enumerate}
  \item $\P(\Omega)=1$;
  \item if $A_1,A_2,\ldots$ are pairwise disjoint, then
  \[
  \P\!\left(\bigcup_i A_i\right)=\sum_i\P(A_i).
  \]
\end{enumerate}
\end{definition}

The collection $\mathcal F$ is called a \emph{sigma-algebra}. For a finite
sample space we usually take $\mathcal F=2^\Omega$, the collection of all
subsets of $\Omega$. The more careful notation becomes important for
continuous spaces.

\section{Building a finite probability space}

Choose nonnegative numbers $w_\omega$ for the individual outcomes such that
$\sum_{\omega\in\Omega}w_\omega=1$, and define
\[
  \P(A)=\sum_{\omega\in A}w_\omega.
\]
The total mass is one; for disjoint events, the sums use disjoint outcomes,
which gives additivity. Since $\Omega$ is finite, only finitely many events
in a disjoint family can be nonempty, so countable additivity also follows.
This construction gives every probability measure on a finite space with
all subsets measurable. The uniform model is the special choice
$w_\omega=1/|\Omega|$ for each outcome.

\begin{example}[Unequal weights]
Let $\Omega=\{a,b,c,d\}$ have masses $1/10,2/10,3/10,4/10$.
For $A=\{a,b\}$ and $B=\{b,d\}$,
\[
\P(A)=\frac3{10},\quad \P(B)=\frac6{10},\quad
\P(A\cap B)=\frac2{10},\quad \P(A\cup B)=\frac7{10}.
\]
Although $A$ contains half the outcomes, it does not have probability
$1/2$. Counting outcomes and adding their probability masses are different
operations unless all outcomes have equal mass.
\end{example}

\begin{proposition}
For events $A,B$,
\[
\P(A^c)=1-\P(A),\qquad
\P(A\cup B)=\P(A)+\P(B)-\P(A\cap B).
\]
If $A\subseteq B$, then $\P(A)\le \P(B)$.
\end{proposition}

Each statement follows by splitting a set into disjoint pieces. This is the basic proof move in probability.

\begin{proof}[The two-event formula]
Write $A\cup B=A\sqcup(B\setminus A)$, where $\sqcup$ denotes disjoint
union. Also $B=(A\cap B)\sqcup(B\setminus A)$. Additivity gives
$\P(A\cup B)=\P(A)+\P(B\setminus A)$ and
$\P(B\setminus A)=\P(B)-\P(A\cap B)$; substitute the second identity
into the first.
\end{proof}

\section{Inclusion--exclusion for probabilities}

For any three events, the same overlap correction gives
\begin{align*}
\P(A\cup B\cup C)
&=\P(A)+\P(B)+\P(C)\\
&\quad-\P(A\cap B)-\P(A\cap C)-\P(B\cap C)
  +\P(A\cap B\cap C).
\end{align*}
No independence or equal-likelihood assumption is needed.
One proof applies the two-event formula to $(A\cup B)\cup C$, then
to $(A\cap C)\cup(B\cap C)=(A\cup B)\cap C$.

\begin{theorem}[Finite probabilistic inclusion--exclusion]
For events $A_1,\ldots,A_r$ in any probability space,
\[
\P\!\left(\bigcup_{i=1}^r A_i\right)
=\sum_{\varnothing\ne J\subseteq\{1,\ldots,r\}}
(-1)^{|J|+1}\P\!\left(\bigcap_{j\in J}A_j\right).
\]
\end{theorem}
\begin{proof}
Partition the sample space by the exact list of events that occur.
There are at most $2^r$ such disjoint measurable regions. A region in
exactly $t\ge1$ events contributes its probability mass with coefficient
$\sum_{k=1}^t(-1)^{k+1}\binom tk=1$ on the right. The region in no event
contributes zero. Finite additivity therefore gives the probability of
the union, even when the original sample space is infinite.
\end{proof}

\begin{example}[Counts become probabilities under a uniform model]
Choose a student uniformly from the 100 students in the activity example
in Chapter 1. Dividing all its counts by 100 gives
\[
\P(A\cup B\cup C)=.39+.29+.25-.13-.11-.09+.05=.65.
\]
The eight disjoint regions have probabilities
$.20,.12,.10,.08,.06,.04,.05,.35$ (the three single-only regions,
three pair-only regions, triple region, and outside region).
They are nonnegative and sum to one, so they also specify a valid
eight-outcome probability space with \emph{unequal} outcome weights.
\end{example}

\begin{warningbox}
Inclusion--exclusion always uses the \emph{actual} intersection
probabilities. Replacing $\P(A\cap B)$ with $\P(A)\P(B)$ requires an
additional independence assumption. The formula itself makes no such
assumption. Likewise, counts divided by a total give probabilities only
for a justified uniform model.
\end{warningbox}

\section{Equally likely is an assumption}

If a finite sample space has $N$ equally likely outcomes and $A$ contains $m$, then $\P(A)=m/N$. But equality of probabilities is part of the model, not a consequence of counting.

\begin{example}
For two independent fair dice, the ordered-pair space has 36 equally likely outcomes. The possible sums $2,3,\ldots,12$ are not equally likely: sum 7 has six representations, while sum 2 has one.
\end{example}

Recording a coarser outcome carries its probabilities along: for a recorded
value $s$, add the masses of all original outcomes producing $s$.
Replacing the resulting weights by equal weights changes the experiment.

\section{The birthday paradox: count the complement}
\label{sec:birthday}

\begin{question}
How large must a class be before a shared birthday is more likely than not?
\end{question}

State the model first. There are $n$ labeled people and $d\ge1$ possible days.
Their birthdays are independent and uniform on those days. Thus
\[
\Omega=\{1,\ldots,d\}^n,\qquad \mathcal F=2^\Omega,
\qquad \P(\{\omega\})=d^{-n}.
\]
For the usual classroom model set $d=365$: ignore leap days and seasonal
variation. Independence and equal daily probabilities are modeling
assumptions, not claims about every real class.

\begin{proposition}[Exact collision probability]
For $0\le n\le d$, the probability $q_d(n)$ of at least one shared
birthday is
\[
q_d(n)=1-\frac{d(d-1)\cdots(d-n+1)}{d^n}
      =1-\prod_{j=0}^{n-1}\left(1-\frac jd\right).
\]
The empty product is one, so $q_d(0)=q_d(1)=0$.
For $n>d$, $q_d(n)=1$ by the pigeonhole principle.
\end{proposition}
\begin{proof}
Count the complement: all birthdays are distinct. There are $d$ choices
for the first birthday, $d-1$ for the second, and so on, yielding $(d)_n$
ordered birthday lists. All $d^n$ lists are equally likely, so divide and
subtract from one. Counting \emph{at least one} matching pair directly is
harder because different matching-pair events overlap.
\end{proof}

\begin{center}
\begin{tabular}{@{}rrrrrr@{}}
\toprule
Class size $n$ & 10 & 20 & 22 & 23 & 30 \\
\midrule
$q_{365}(n)$ & .11695 & .41144 & .47570 & .50730 & .70632 \\
\bottomrule
\end{tabular}
\end{center}
Adding a person cannot destroy an existing match, so the probability is
nondecreasing. Since $q_{365}(22)<1/2<q_{365}(23)$, the first class size
above $50\%$ is $23$.

\begin{checkpoint}[title={Any pair versus one particular person}]
In a class of $n\ge1$ people, the probability that someone else shares one
\emph{designated} person's birthday is
\[
1-\left(1-\frac1d\right)^{n-1}.
\]
For $d=365,n=23$, this is about $.05857$, much smaller than $.50730$.
The designated-person question compares that person with 22 others;
the any-pair question considers $\binom{23}{2}=253$ possible pairs.
\end{checkpoint}

\section{Birthday matches meet inclusion--exclusion}

For three people, let $E_{12},E_{13},E_{23}$ be the events that the
indicated people have the same birthday. Each event has probability $1/d$.
Any two of these events force all three birthdays to agree, so each
pairwise intersection has probability $1/d^2$. The triple intersection
is the \emph{same event}, with probability $1/d^2$, not $1/d^3$.
Therefore inclusion--exclusion gives
\[
q_d(3)=\frac3d-\frac3{d^2}+\frac1{d^2}
       =\frac3d-\frac2{d^2},
\]
agreeing with $1-(1-1/d)(1-2/d)$. This illustrates why treating all
matching-pair events as independent gives a wrong answer.

More generally there are $\binom n2$ matching-pair events, each with
probability $1/d$. The union bound gives the useful check
\[
q_d(n)\le\min\!\left(1,\frac{\binom n2}{d}\right).
\]
The sum $\binom n2/d$ counts overlapping pair events; it is not itself
the collision probability. Later, linearity of expectation will identify
it as the expected number of matching pairs.

\begin{labbox}[title={The birthday experiment}]
Use the website to compare the exact product, the designated-person
probability, and the union bound. Predict the $50\%$ threshold, then
simulate classes. Change the calendar from 365 days to 30 or 10 days.
Separate sampling fluctuations from changes in the underlying model.
\end{labbox}

\section{Continuity and geometry}

Choose a point uniformly from $[0,1]$. Every individual point has probability zero, yet the entire interval has probability one. Probabilities attach to regions: $\P([a,b])=b-a$. Zero probability does not mean impossible; it means negligible under the model.

\begin{center}
\begin{tikzpicture}[x=10cm,y=1cm]
  \draw[indigo,thick] (0,0)--(1,0);
  \draw[mint,line width=5pt] (.24,0)--(.71,0);
  \foreach \x/\lab in {0/0,.24/a,.71/b,1/1}{\draw (\x,-.10)--(\x,.10) node[below=4pt]{\small $\lab$};}
  \node[above=6pt,mint!70!black] at (.475,0) {$\P([a,b])=b-a$};
\end{tikzpicture}
\end{center}

\section{Simulation and the model--world gap}

A probability model is a deliberate simplification. Simulation samples from the model; data come from the world. Agreement between simulation and data supports a model, but never proves it uniquely correct.

\begin{theorem}[Union bound]
For any events $A_1,\ldots,A_n$,
\[
\P\!\left(\bigcup_{i=1}^n A_i\right)\le \sum_{i=1}^n\P(A_i).
\]
\end{theorem}

The events need not be disjoint. Double counting makes the right side larger, which is harmless for an upper bound.

\begin{exercisebox}
\begin{enumerate}
  \item Derive $\P(\varnothing)=0$ from the axioms. \exlevel{Reasoning}
  \item Two independent fair dice are rolled. Find the probability that the maximum is 4. \exlevel{Fluency}
  \item Prove the union bound for three events. \exlevel{Reasoning}
  \item A point is chosen uniformly from the unit square. Find the probability that $x+y<1$. \exlevel{Fluency}
  \item Give an event with probability zero that is not logically impossible. \exlevel{Reasoning}
  \item Estimate by simulation the probability that two random points in $[0,1]$ are within $0.1$ of each other; then derive it geometrically. \exlevel{Exploration}
  \item On $\Omega=\{1,2,3,4\}$, set $\P(\{i\})=i/10$. Compute the probability of an even outcome and compare it with the fraction of outcomes that are even. \exlevel{Fluency}
  \item Prove three-event probabilistic inclusion--exclusion using only the two-event formula. \exlevel{Reasoning}
  \item With independent uniform birthdays on 30 days, find the first class size whose collision probability exceeds $1/2$. \exlevel{Fluency}
  \item For four people and a five-day calendar, count the no-collision lists and calculate the collision probability. \exlevel{Fluency}
  \item Explain why $\binom n2/d$ can exceed one without contradicting the probability axioms. \exlevel{Reasoning}
\end{enumerate}
\end{exercisebox}

\chapter{Information changes the space}
\chaptertag{Conditional probability, independence, and Bayes' rule}

\begin{question}
When evidence arrives, which probabilities should change---and by how much?
\end{question}

\begin{definition}
If $\P(B)>0$, the conditional probability of $A$ given $B$ is
\[
\P(A\mid B)=\frac{\P(A\cap B)}{\P(B)}.
\]
\end{definition}

Conditioning restricts attention to $B$ and renormalizes so that $B$ becomes the new whole.

\begin{center}
\begin{tikzpicture}[scale=.85]
  \fill[indigo!8] (-3,-1.5) rectangle (3,1.6);
  \fill[amber!40] (-.4,0) circle (1.2);
  \fill[mint!35] (.9,0) circle (1.2);
  \draw[amber!80!black,thick] (-.4,0) circle (1.2) node[left=8mm] {$A$};
  \draw[mint!80!black,thick] (.9,0) circle (1.2) node[right=8mm] {$B$};
  \node[below=16mm] at (.9,0) {Within $B$, measure the part also in $A$};
\end{tikzpicture}
\end{center}

\section{Multiplication and trees}

Rearranging the definition gives
\[
\P(A\cap B)=\P(B)\P(A\mid B)=\P(A)\P(B\mid A).
\]
Along a probability tree, multiply down branches and add across disjoint terminal paths.

\section{Independence}

Events $A$ and $B$ are independent when
\[
\P(A\cap B)=\P(A)\P(B).
\]
When probabilities are positive, this is equivalent to $\P(A\mid B)=\P(A)$. Disjointness is almost the opposite: two nontrivial disjoint events cannot be independent.

\begin{warningbox}
Pairwise independence does not imply mutual independence. Three events can be independent in every pair while the third is determined by the first two.
\end{warningbox}

\section{Total probability and Bayes}

If $B_1,\ldots,B_k$ partition $\Omega$, then
\[
\P(A)=\sum_{i=1}^k \P(A\mid B_i)\P(B_i).
\]
Bayes' rule reverses the condition:
\[
\P(B_j\mid A)=
\frac{\P(A\mid B_j)\P(B_j)}{\sum_i \P(A\mid B_i)\P(B_i)}.
\]

\begin{example}[A diagnostic test]
Suppose prevalence is $1\%$, sensitivity is $95\%$, and specificity is $90\%$. Among 10,000 people, expect 100 cases, of which 95 test positive. Among 9,900 non-cases, 990 test positive. Thus
\[
\P(\text{case}\mid +)=\frac{95}{95+990}\approx 0.0876.
\]
A highly sensitive test can have a modest positive predictive value when the base rate is low.
\end{example}

\begin{labbox}
Represent the test example with 10,000 dots rather than formulas. Change prevalence while holding sensitivity and specificity fixed. Plot the positive predictive value; locate the prevalence at which a positive result becomes more likely correct than incorrect.
\end{labbox}

\section{A preview: symmetry is not independence}

Choose once between coins with heads probabilities $1/4$ and $3/4$, giving each coin probability $1/2$, and then repeatedly toss the chosen coin. Conditional on the choice, the tosses are independent. Without that information, if $H_i$ means heads on toss $i$, then
\[
\P(H_1)=\P(H_2)=\frac12,\qquad
\P(H_1\cap H_2)=\frac12\left(\frac14\right)^2
+\frac12\left(\frac34\right)^2=\frac5{16}\ne\frac14.
\]
The tosses are dependent, yet any reordering of a finite string has the same probability as the original string. This symmetry is called \emph{exchangeability}. Chapter~\ref{ch:exchangeability} returns to the example and proves de Finetti's surprising converse for infinite zero--one sequences: exchangeability gives a representation by a shared random bias and conditionally independent trials.

\begin{exercisebox}
\begin{enumerate}
  \item Draw a tree for two draws without replacement from an urn with three red and two blue balls. \exlevel{Fluency}
  \item Prove that if $A$ and $B$ are independent, then $A^c$ and $B$ are independent. \exlevel{Reasoning}
  \item Construct three pairwise independent events that are not mutually independent. \exlevel{Reasoning}
  \item Solve the diagnostic-test example with prevalence $5\%$. \exlevel{Fluency}
  \item A randomly selected coin is fair with probability $.8$ and double-headed with probability $.2$. It lands heads three times. Find the posterior probability it is double-headed. \exlevel{Reasoning}
  \item Explain precisely why changing the information protocol can change the ``two children'' answer. \exlevel{Exploration}
  \item Simulate Monty Hall and explain the result using conditioning rather than frequency alone. \exlevel{Exploration}
\end{enumerate}
\end{exercisebox}

\chapter{Numbers attached to outcomes}
\chaptertag{Random variables, distributions, expectation, and moments}

\begin{question}
How can one function compress a complicated experiment into a numerical quantity of interest?
\end{question}

\begin{definition}
A random variable is a function $X:\Omega\to\R$. Its cumulative distribution function is
\[
F_X(x)=\P(X\le x).
\]
\end{definition}

Every CDF is nondecreasing, right-continuous, approaches 0 to the left and 1 to the right. For a discrete variable, jumps encode probabilities. For a continuous variable with density $f$, areas encode probabilities:
\[
F_X(x)=\int_{-\infty}^x f(t)\,dt,
\qquad
\P(a<X\le b)=\int_a^b f(t)\,dt.
\]

\section{Expectation}

For discrete $X$ with values $x_i$,
\[
\E[X]=\sum_i x_i\P(X=x_i).
\]
For continuous $X$ with density $f$,
\[
\E[X]=\int_{-\infty}^{\infty}x f(x)\,dx.
\]
Expectation is a center of mass, not necessarily a possible value.

\begin{theorem}[LOTUS]
For a suitable function $g$,
\[
\E[g(X)]=\sum_x g(x)\P(X=x)
\quad\text{or}\quad
\E[g(X)]=\int g(x)f_X(x)\,dx.
\]
\end{theorem}

\begin{theorem}[Linearity]
For constants $a,b$ and integrable random variables $X,Y$,
\[
\E[aX+bY]=a\E[X]+b\E[Y],
\]
whether or not $X$ and $Y$ are independent.
\end{theorem}

\section{Variance and moments}

\[
\Var(X)=\E[(X-\E X)^2]=\E[X^2]-(\E X)^2.
\]
Standard deviation has the same units as $X$. The $k$th raw moment is $\E[X^k]$; the $k$th central moment is $\E[(X-\E X)^k]$. The moment-generating function, when finite near zero, is $M_X(t)=\E[e^{tX}]$ and satisfies $M_X^{(k)}(0)=\E[X^k]$.

\begin{proposition}
For constants $a,b$,
\[
\Var(aX+b)=a^2\Var(X).
\]
\end{proposition}

\section{Indicator variables}

For an event $A$, let $\ind_A=1$ on $A$ and $0$ otherwise. Then $\E[\ind_A]=\P(A)$. If $N=\sum_i\ind_{A_i}$ counts occurrences, linearity gives
\[
\E[N]=\sum_i\P(A_i)
\]
without independence.

\begin{theorem}[Markov and Chebyshev]
If $X\ge0$, then $\P(X\ge a)\le \E[X]/a$. Consequently,
\[
\P(|X-\E X|\ge k\sigma)\le \frac1{k^2}.
\]
\end{theorem}

\begin{exercisebox}
\begin{enumerate}
  \item Find the CDF of the number of heads in two fair flips. \exlevel{Fluency}
  \item A density is proportional to $x$ on $[0,2]$. Normalize it and find its mean. \exlevel{Fluency}
  \item Prove $\Var(X)=\E[X^2]-(\E X)^2$. \exlevel{Reasoning}
  \item Find the expected number of fixed points in a random permutation of $n$ objects. \exlevel{Reasoning}
  \item Show that a fair lottery can still have a large variance. Interpret. \exlevel{Exploration}
  \item Use Chebyshev's inequality to bound the chance that a variable lies at least three standard deviations from its mean. Compare with the normal distribution. \exlevel{Reasoning}
  \item Construct two different distributions with the same mean and variance. \exlevel{Exploration}
\end{enumerate}
\end{exercisebox}

\chapter{A library of chance mechanisms}
\label{ch:standard-distributions}
\chaptertag{Discrete distributions first; continuous models later}

\begin{question}
Which probability laws recur because the same mechanisms recur?
\end{question}

For a discrete variable, its distribution lists the possible values and
their probabilities.
The coin counts, dice, and hat counts of Chapter 4 already have
distributions. A \emph{named family} collects experiments whose probability
tables have the same form. Learn the experiment before memorizing the name.

\begin{bigidea}[title={A route through the next topics}]
Start with finite distributions: a uniform choice, one success indicator,
a count of independent successes, and sampling without replacement.
For each, identify the possible values, calculate their probabilities,
then find the mean and variance.
Geometric waiting times and Poisson counts come next as countably infinite
examples. The coupon collector then brings waiting times, expectation,
variance, and inclusion--exclusion together. The continuous section is
for later, after the calculus refresher.
\end{bigidea}

\section{A uniform choice from a finite list}
\label{sec:discrete-families}

Choose one of $1,\ldots,m$ uniformly, where $m\ge1$. Then
$\P(U=k)=1/m$ for $k=1,\ldots,m$, and zero for other values.
The usual sum and sum-of-squares identities give
\[
 \E[U]=\frac{m+1}{2},\qquad
 \E[U^2]=\frac{(m+1)(2m+1)}6,\qquad
 \Var(U)=\frac{m^2-1}{12}.
\]
For example, choosing a label uniformly from $1,\ldots,8$ gives
mean $9/2$ and variance $21/4$.
The case $m=1$ is constant, with variance zero.

A uniform distribution on any finite set $\{u_1,\ldots,u_m\}$ averages
those values with weights $1/m$. The formulas above are specifically for
the consecutive integers $1,\ldots,m$.
Uniform outcomes also need not produce uniform recorded values:
the sum of two fair dice is not uniform on $2,\ldots,12$.

\section{Bernoulli: one event, recorded as zero or one}
\label{sec:bernoulli}

If an event has probability $p$, its indicator $I$ has distribution
\[
 \P(I=1)=p,\qquad \P(I=0)=1-p,\qquad 0\le p\le1.
\]
We call this the \emph{Bernoulli distribution} with parameter $p$.
There are only two values, so directly
\[
 \E[I]=p,\qquad \E[I^2]=p,\qquad \Var(I)=p(1-p).
\]
The parameter $p$ is a probability, the value of $I$ is zero or one,
and its expectation can lie between these values.
At $p=0$ or $1$ the result is certain and the variance vanishes.

\begin{example}
Let $I$ record whether a fair die shows a six. Then $p=1/6$,
$\E[I]=1/6$, and $\Var(I)=5/36$.
A success may mean a correct answer, a defective item, or a birthday
match for a specified pair; a \emph{success} need not mean something desirable.
\end{example}

\section{Binomial: count successes in a fixed number of trials}
\label{sec:binomial}

Perform $n$ independent trials, each with success probability $p$,
and let $X$ count the successes. For $0<p<1$ and $k=0,\ldots,n$,
\[
 \P(X=k)=\binom nkp^k(1-p)^{n-k}.
\]
A specified success/failure string has probability $p^k(1-p)^{n-k}$;
choose its $k$ success positions in $\binom nk$ ways.
The binomial theorem verifies that the masses sum to one.
We write $X\sim\Bin(n,p)$.
If $p=0$, $X=0$; if $p=1$, $X=n$; if $n=0$, $X=0$.

Write $X=I_1+\cdots+I_n$. Linearity and the independent-variance rule give
\[
 \E[X]=np,\qquad \Var(X)=np(1-p).
\]
This calculation avoids summing $k^2$ against the binomial probabilities.

\begin{example}[Five trials]
For $X\sim\Bin(5,1/3)$,
\[
 \P(X=3)=\binom53(1/3)^3(2/3)^2=\frac{40}{243},
 \qquad \E[X]=\frac53,\qquad \Var(X)=\frac{10}{9}.
\]
For a threshold it can be easier to use a complement:
$\P(X\ge2)=1-\P(X=0)-\P(X=1)$.
\end{example}

\begin{warningbox}
A count is binomial only when its experiment satisfies the assumptions:
a fixed number of trials, two outcomes per trial, a common success
probability, and independence.
Recording the same toss ten times gives a count of either zero or ten,
not a binomial count of ten independent tosses.
Sampling from a finite population without replacement also requires
a different model.
\end{warningbox}

\section{Hypergeometric: sample without replacement}
\label{sec:hypergeometric}

A population contains $K$ marked objects among $N$ distinct objects.
Choose a subset of $m$ objects uniformly, and let $X$ count marked objects
in the sample. There are $\binom Nm$ equally likely samples.
To obtain exactly $k$ marked objects, choose those $k$ from the marked
group and the other $m-k$ from the unmarked group:
\[
 \P(X=k)=\frac{\binom Kk\binom{N-K}{m-k}}{\binom Nm}.
\]
Here $N\ge1$, $0\le K,m\le N$, and the possible values satisfy
$\max(0,m-(N-K))\le k\le\min(m,K)$.
The formula is normalized because these disjoint cases partition all
$m$-element samples. This is the \emph{hypergeometric distribution}.

\begin{example}[Two draws from five balls]
An urn contains three red and two blue balls. For a uniform sample of
two balls without replacement, the number $X$ of red balls has table
\[
\begin{array}{c|ccc}
k&0&1&2\\\hline
\P(X=k)&1/10&6/10&3/10
\end{array}
\]
Thus $\E[X]=6/5$, $\E[X^2]=9/5$, and $\Var(X)=9/25$.
With replacement the count would instead be $\Bin(2,3/5)$, with the
same mean but variance $12/25$.
\end{example}

For the general moments, regard the sample as the first $m$ positions
in a uniform random ordering of the population.
Put $p=K/N$ and let $I_j$ indicate a marked object in position $j$.
Every position has marked probability $p$, so $\E[X]=mp$.
For $N\ge2$ and $i\ne j$,
\[
 \E[I_iI_j]=\frac{K(K-1)}{N(N-1)},\qquad
 \E[I_iI_j]-p^2=-\frac{p(1-p)}{N-1}.
\]
The variance-of-a-sum expansion from Chapter 4 gives
\[
 \Var(X)=mp(1-p)-m(m-1)\frac{p(1-p)}{N-1}
        =mp(1-p)\frac{N-m}{N-1}.
\]
For $1\le m\le N$, the last factor is at most one: sampling without
replacement has no more variance than the corresponding binomial count.
When $m=N$, the count is the fixed number $K$, so its variance is zero.
For a one-object population the count is constant as well.

\section{For later: geometric waiting times}
\label{sec:geometric}

Now allow a countably infinite list of possible values.
In independent trials with fixed success probability $0<p\le1$,
let $T$ count trials up to and including the first success. Then
\[
 \P(T=k)=(1-p)^{k-1}p,\qquad k=1,2,\ldots.
\]
For $0<p<1$, the first $k-1$ trials must fail and the $k$th succeed.
The geometric series sums the masses to one.
For $p=1$, $T=1$ with certainty. If $p=0$, there is no finite waiting
time to a success, so it is excluded from this distribution.

The event $T>r$ means $r$ initial failures, hence
$\P(T>r)=(1-p)^r$ for nonnegative integers $r$.
For $r=0$ the probability is one, including at $p=1$.
For $0<p<1$ and integers $s,t\ge0$,
\[
 \P(T>s+t\mid T>s)
 =\frac{(1-p)^{s+t}}{(1-p)^s}
 =(1-p)^t=\P(T>t).
\]
This is \emph{memorylessness}: after any number of failures, the
distribution of the additional wait is unchanged.
For $p=1$, only conditioning events of positive probability can be used.

\subsection*{Mean and variance from tail sums}
A positive integer $T$ satisfies
$T=\sum_{j\ge1}\ind_{\{T\ge j\}}$ and
$T^2=\sum_{j\ge1}(2j-1)\ind_{\{T\ge j\}}$.
First truncate these nonnegative sums; letting the truncation grow gives
\[
 \E[T]=\sum_{j\ge1}\P(T\ge j)=\sum_{j\ge1}(1-p)^{j-1}=\frac1p.
\]
The geometric-series calculation
$\sum_{j\ge1}j r^{j-1}=1/(1-r)^2$, reviewed in
Section~\ref{sec:calc-series}, similarly gives
\[
 \E[T^2]=\frac2{p^2}-\frac1p=\frac{2-p}{p^2},
 \qquad \Var(T)=\frac{1-p}{p^2}.
\]
These sums converge for $p>0$; at $p=1$ the values follow directly.
Waiting for a six on independent fair rolls gives mean $6$,
variance $30$, and $\P(T>5)=(5/6)^5$.

\begin{warningbox}
Some books count \emph{failures before} the first success.
That variable is $T-1$, supported on $0,1,\ldots$, with mean
$(1-p)/p$ and the same variance. State which convention you use.
\end{warningbox}

\section{For later: Poisson counts}
\label{sec:poisson}

For $\lambda>0$, the \emph{Poisson distribution} assigns
\[
 \P(X=k)=e^{-\lambda}\frac{\lambda^k}{k!},
 \qquad k=0,1,2,\ldots.
\]
The exponential series makes the total mass
$e^{-\lambda}\sum_{k\ge0}\lambda^k/k!=1$.
The limiting case $\lambda=0$ places all mass at zero.
Shifting indices in absolutely convergent sums gives
\[
 \E[X]=e^{-\lambda}\lambda\sum_{j\ge0}\frac{\lambda^j}{j!}
 =\lambda,\qquad
 \E[X(X-1)]=\lambda^2.
\]
Consequently $\E[X^2]=\lambda^2+\lambda$ and $\Var(X)=\lambda$.
Factorial moments make this calculation short, just as they did for hats.

The law arises from many independent rare opportunities. For fixed
$\lambda>0$ and fixed integer $k\ge0$,
\[
 \binom nk(\lambda/n)^k(1-\lambda/n)^{n-k}
 \longrightarrow e^{-\lambda}\frac{\lambda^k}{k!}.
\]
Indeed, $(n)_k/n^k\to1$ and $(1-\lambda/n)^n\to e^{-\lambda}$.
This explains a Poisson approximation to $\Bin(n,p)$ when $p$ is small
and $np$ is near $\lambda$; it does not make every count Poisson.
In particular, the model's equality of mean and variance is a consequence
to check, not a rule for all counts.

For $X$ Poisson with mean $3$, $\P(X\ge2)=1-4e^{-3}$.
The number of fixed points of a large uniform permutation approaches
Poisson with mean one, as derived in
Section~\ref{sec:calc-derangements}.
For a finite permutation this is an approximation; the exact law is
$\P(N=k)=\binom nkD_{n-k}/n!$.

\section{Coupon collection: waiting to see every type}
\label{sec:coupon-collector}

An album has $n\ge1$ types of coupon. Each purchase gives one coupon,
independently and uniformly chosen from the $n$ types; duplicates are
possible. Let $T_n$ be the number of purchases needed to see every type,
starting with an empty album. Equivalently, how many independent rolls
of a fair die are needed to see all six faces?

This reverses the birthday question. There we stopped at the first
repeat; here we keep going until nothing is missing.
With sampling without replacement, completion would take exactly $n$
draws. Repeated types are what make this problem interesting.

\subsection*{Break the wait into stages}
After $k$ distinct types have appeared, let $W_k$ count additional draws
up to and including the next new type, for $k=0,\ldots,n-1$.
There are $n-k$ unseen types. Until one appears, the chance of a new type
on each draw stays at
\[
 p_k=\frac{n-k}{n},\qquad
 W_k\sim\text{Geometric}(p_k),\qquad
 \E[W_k]=\frac{n}{n-k}.
\]
In particular $W_0=1$: the first draw is always new.
The collection time is the sum of these stages,
\[
 T_n=W_0+W_1+\cdots+W_{n-1}.
\]
Linearity of expectation now does the work. Define the \emph{harmonic
numbers} by $H_0=0$ and $H_n=\sum_{j=1}^n1/j$. Then
\begin{equation}
 \boxed{\E[T_n]=\frac nn+\frac n{n-1}+\cdots+\frac n1=nH_n.}
 \label{eq:coupon-mean}
\end{equation}
Every stage has finite expectation, so completion occurs with probability
one, although there is no fixed upper limit on the wait when $n\ge2$.
The last missing type alone takes $n$ further draws on average.

\begin{example}[Every face of a die]
For six equally likely types, the successive expected waits are
\[
 1,\quad\frac65,\quad\frac64,\quad\frac63,\quad\frac62,\quad6.
\]
Their sum is $6H_6=14.7$ rolls. Seeing the first five types takes $8.7$
rolls on average, and the last one takes another six.
The mean is an average over repetitions of the entire experiment,
not a deadline by which every collection will finish.
\end{example}

For a sense of scale, comparing rectangles with the area under $1/x$
gives, for $n\ge1$,
\[
 \log(n+1)=\int_1^{n+1}\frac{dx}{x}
 \le H_n\le1+\int_1^n\frac{dx}{x}=1+\log n.
\]
Thus $\E[T_n]$ grows like $n\log n$, rather than a fixed multiple of $n$.
Here and below, $\log$ is the natural logarithm. The integral comparison
can wait until Appendix~\ref{app:calculus}; the exact mean above uses
only geometric expectations and a finite sum.

\subsection*{Variance: justify independence before adding}
The expectation calculation did not require independence of the stages.
For the variance, we do need it. After any specified completed history,
future draws are still independent and uniform. The conditional law of
$W_k$ is always geometric with parameter $(n-k)/n$, regardless of which
types have appeared or how long earlier stages took. Applying the
multiplication rule to these conditional probabilities shows that
$W_0,\ldots,W_{n-1}$ are mutually independent.

The geometric variance formula therefore gives
\begin{align*}
 \Var(T_n)
 &=\sum_{k=0}^{n-1}\frac{1-p_k}{p_k^2}
   =\sum_{j=1}^n\frac{1-j/n}{(j/n)^2}\\
 &=\boxed{n^2\sum_{j=1}^n\frac1{j^2}-nH_n.}
\end{align*}
For a die, this is $38.99$ rolls squared, giving standard deviation
about $6.24$ rolls. For $n=1$, $T_1=1$ and the variance is zero.
For $n\ge2$, the entire collection time is not itself geometric:
the success probability changes each time a new type is found.

\subsection*{A deadline: count the types still missing}
Fix an integer $m\ge1$. Let $M_m$ be the number of types missing after
$m$ draws, and let $A_i$ be the event that type $i$ is missing.
Then
\[
 M_m=\sum_{i=1}^n\ind_{A_i},\qquad
 \P(A_i)=\left(1-\frac1n\right)^m,\qquad
 \E[M_m]=n\left(1-\frac1n\right)^m.
\]
These missing-type indicators are generally dependent. Linearity still
applies, and the union bound (equivalently, Markov at one) gives
\begin{equation}
 \P(T_n>m)=\P(M_m\ge1)
 \le\min\left\{1,n\left(1-\frac1n\right)^m\right\}
 \le\min\{1,ne^{-m/n}\}.
 \label{eq:coupon-tail-bound}
\end{equation}
We used $1-u\le e^{-u}$. Consequently, for $0<\delta<1$, any integer
$m\ge n(\log n+\log(1/\delta))$ guarantees completion probability
at least $1-\delta$. This is a sufficient bound, not the exact cutoff.

Inclusion--exclusion gives an exact answer: a specified set of $j$ types
is missing with probability $(1-j/n)^m$, so
\begin{equation}
 \P(T_n\le m)=\sum_{j=0}^n(-1)^j\binom nj
                         \left(1-\frac jn\right)^m,\qquad m\ge1.
 \label{eq:coupon-cdf}
\end{equation}
At $m=0$ the completion probability is zero and $M_0=n$.
For a die, the formula gives $\P(T_6\le15)\approx.6442$:
a deadline close to the mean is far from a guarantee.

\subsection*{Bonus: a Poisson number of draws}
There is also a direct connection to the Poisson distribution, entirely
within discrete probability. Choose $K\sim\Pois(t)$, where $t>0$,
independently of the coupon draws, and collect exactly $K$ coupons.
Let $C_i$ be the number of coupons of type $i$.
For nonnegative integers $c_1,\ldots,c_n$ with sum $s$, first condition
on $K=s$, then count the orderings of those coupon labels:
\begin{align*}
 \P(C_1=c_1,\ldots,C_n=c_n)
 &=\frac{e^{-t}t^s}{s!}\,
   \frac{s!}{c_1!\cdots c_n!}\left(\frac1n\right)^s\\
 &=\prod_{i=1}^n e^{-t/n}\frac{(t/n)^{c_i}}{c_i!}.
\end{align*}
The joint probability factors into Poisson probabilities. Thus the
$C_i$ are independent Poisson variables of mean $t/n$.
Each type is missing with probability $e^{-t/n}$, independently of
the other types. Hence
\[
 M_K\sim\Bin(n,e^{-t/n}),\qquad
 \P(T_n\le K)=\P(M_K=0)=\left(1-e^{-t/n}\right)^n.
\]
This trick is called \emph{Poissonization}. With a fixed number $m$ of
draws, the counts sum to $m$ and are dependent when $n\ge2$ and $m\ge1$.
Randomizing their total in this particular way makes them independent.
The product formula is exact for a Poisson number of draws; it is not
the fixed-$m$ answer in~\eqref{eq:coupon-cdf}.
No continuous waiting-time model is needed for this argument.

\section{Practice: choose and use a discrete model}
\label{sec:discrete-distribution-exercises}

\begin{exercisebox}
\begin{enumerate}
 \item For $X\sim\Bin(5,.3)$, find $\P(X\ge2)$, its mean, and variance.
 \exlevel{Fluency}
 \item Prove the geometric memoryless property, stating the allowed
 values and the condition that the conditioning event has positive
 probability. \exlevel{Reasoning}
 \item Approximate $\Bin(1000,.002)$ by a Poisson law and compare
 the exact and approximate $\P(X=0)$. \exlevel{Exploration}
 \item Choose a label uniformly from $1,\ldots,8$. Find its mean,
 variance, and probability of exceeding its mean. \exlevel{Fluency}
 \item An event has probability $.2$. Find the mean and variance
 of its indicator. Compare counting successes in ten independent
 repetitions with recording one indicator ten times. \exlevel{Reasoning}
 \item Draw two balls without replacement from three red and two blue
 balls. Derive the complete red-count distribution, mean, and variance.
 Repeat with replacement. \exlevel{Fluency}
 \item For each discrete family, describe one experiment that fits
 and one that violates a stated modeling assumption. \exlevel{Exploration}
 \item Roll a fair die independently until the first six. Find the
 mean and variance of the number of rolls and the probability it exceeds
 five. What is the mean number of failures? \exlevel{Fluency}
 \item For a Poisson variable of mean three, find $\P(X\ge2)$ and
 $\E[X(X-1)]$. Recover its variance. \exlevel{Fluency}
 \item A population has six marked objects among twenty. Take four
 draws. Compare the count's mean and variance with and without replacement.
 State why the probability models differ. \exlevel{Reasoning}
 \item Roll a fair die until all faces have appeared. Find the mean,
 variance, and expected additional wait when exactly two faces are still
 missing. \exlevel{Fluency}
 \item Let $T_{n,r}$ count draws until $r$ distinct types have appeared
 in the uniform coupon model. Show that
 $\E[T_{n,r}]=n(H_n-H_{n-r})$ for $1\le r\le n$.
 What is the expected remaining wait for a complete collection?
 \exlevel{Reasoning}
 \item For 100 equally likely coupon types, use the exponential bound
 in~\eqref{eq:coupon-tail-bound} to choose a whole number of draws
 guaranteeing at least $95\%$ completion probability. Explain why
 $100H_{100}$ is not such a guarantee. \exlevel{Reasoning}
 \item Independently choose $K\sim\Pois(12)$, then roll a fair die $K$ times.
 Find the distribution of the number of unseen faces and the probability
 of seeing every face. Compare with exactly 12 rolls. \exlevel{Exploration}
\end{enumerate}
\end{exercisebox}

\section{For later: continuous families}
\label{sec:continuous-families}

This section uses densities and integrals from
Section~\ref{sec:rv-continuous} and Appendix~\ref{app:calculus}.
It can be left until after the discrete expectation and variance material.

For $a<b$, Uniform$(a,b)$ assigns constant density $1/(b-a)$ to an interval. For $\lambda>0$, Exponential$(\lambda)$ has density $\lambda e^{-\lambda x}$ on $x\ge0$ and is the continuous memoryless law.

The normal density with mean $\mu$ and variance $\sigma^2>0$ is
\[
f(x)=\frac1{\sigma\sqrt{2\pi}}
\exp\!\left(-\frac{(x-\mu)^2}{2\sigma^2}\right).
\]
Standardization sends $X\sim\Normal(\mu,\sigma^2)$ to
\[
Z=\frac{X-\mu}{\sigma}\sim\Normal(0,1).
\]

\begin{center}
\begin{tikzpicture}[x=1.2cm,y=5cm]
  \draw[->,indigo] (-3.5,0)--(3.6,0) node[right] {$z$};
  \draw[mint,very thick,domain=-3.4:3.4,samples=80,smooth] plot (\x,{0.3989*exp(-\x*\x/2)});
  \fill[amber!35,domain=-1:1,samples=40] plot (\x,{0.3989*exp(-\x*\x/2)}) -- (1,0)--(-1,0)--cycle;
  \draw (-1,.012)--(-1,-.012) node[below] {$-1$};
  \draw (1,.012)--(1,-.012) node[below] {$1$};
  \node at (0,.19) {about $68\%$};
\end{tikzpicture}
\end{center}

For $\alpha,\beta>0$, the beta family on $[0,1]$ has density proportional to $x^{\alpha-1}(1-x)^{\beta-1}$. It is useful for proportions and becomes a natural prior for a Bernoulli probability.

\needspace{8\baselineskip}
\section{Choose by mechanism, not by silhouette}

A bell-shaped histogram does not prove normality. Ask what generated the
observations, what values are possible, whether trials are independent,
whether rates are stable, and what approximation error is acceptable.
A named distribution is useful because its assumptions describe an
experiment and its formulas then answer questions about that experiment.

\begin{exercisebox}[title={Later practice: continuous families}]
\begin{enumerate}[start=15]
 \item Find the median of an exponential distribution. \exlevel{Fluency}
 \item If $X\sim\Normal(10,9)$, express $\P(8<X<13)$ using the
 standard normal CDF. \exlevel{Fluency}
 \item Derive the mean and variance of Uniform$(a,b)$ for $a<b$.
 \exlevel{Reasoning}
\end{enumerate}
\end{exercisebox}

\chapter{Variables together}
\chaptertag{Joint, marginal, and conditional distributions}

\begin{question}
How do we describe dependence without confusing association, causation, and prediction?
\end{question}

For discrete variables, a joint mass function $p(x,y)=\P(X=x,Y=y)$ is a table whose entries sum to one. Marginals are obtained by summing:
\[
p_X(x)=\sum_y p(x,y),\qquad p_Y(y)=\sum_x p(x,y).
\]
For a continuous joint density $f(x,y)$, replace sums with integrals.

\begin{center}
\begin{tikzpicture}[scale=.75]
  \foreach \i in {0,...,4}{\foreach \j in {0,...,4}{
    \pgfmathsetmacro{\op}{15+12*\i+8*\j}
    \fill[violet!\op] (\i,\j) rectangle +(1,1);
    \draw[white] (\i,\j) rectangle +(1,1);
  }}
  \draw[->,indigo,thick] (0,-.25)--(5.5,-.25) node[right] {$x$};
  \draw[->,indigo,thick] (-.25,0)--(-.25,5.5) node[above] {$y$};
  \node[right] at (5.6,3.8) {color = joint probability};
\end{tikzpicture}
\end{center}

\section{Conditional laws and independence}

For $p_Y(y)>0$,
\[
p_{X\mid Y}(x\mid y)=\frac{p(x,y)}{p_Y(y)}.
\]
Variables are independent when their joint law factors:
\[
p(x,y)=p_X(x)p_Y(y)
\]
or $f(x,y)=f_X(x)f_Y(y)$ throughout the support.

\section{Covariance and correlation}

\[
\Cov(X,Y)=\E[(X-\E X)(Y-\E Y)]
=\E[XY]-\E[X]\E[Y].
\]
Correlation is the dimensionless quantity
\[
\rho_{X,Y}=\frac{\Cov(X,Y)}{\sigma_X\sigma_Y}.
\]
It lies between $-1$ and $1$. Independence implies zero covariance when moments exist, but zero covariance does not generally imply independence.

\begin{theorem}
\[
\Var(X+Y)=\Var(X)+\Var(Y)+2\Cov(X,Y).
\]
For independent variables, the covariance term vanishes.
\end{theorem}

\section{Conditional expectation}

$\E[X\mid Y]$ is itself a random variable: for each observed value of $Y$, it gives the mean of the corresponding conditional distribution. The tower property says
\[
\E[\E[X\mid Y]]=\E[X].
\]
The conditional mean is also the best mean-square predictor of $X$ using only $Y$.

\begin{warningbox}
Correlation measures linear association. A perfect U-shaped relationship can have correlation zero. And neither correlation nor conditioning alone establishes causation.
\end{warningbox}

\begin{exercisebox}
\begin{enumerate}
  \item Given a $2\times3$ joint probability table, compute both marginals and one conditional distribution. \exlevel{Fluency}
  \item Prove that independence implies $\E[XY]=\E[X]\E[Y]$. \exlevel{Reasoning}
  \item Let $X$ be uniform on $\{-1,0,1\}$ and $Y=X^2$. Show $\Cov(X,Y)=0$ although $X,Y$ are dependent. \exlevel{Reasoning}
  \item Derive the variance-of-a-sum identity. \exlevel{Reasoning}
  \item If $X,Y$ are independent Poisson variables, use generating functions to show $X+Y$ is Poisson. \exlevel{Reasoning}
  \item Create three scatterplots with correlations near $-1$, $0$, and $1$, then create a nonlinear plot with correlation near zero. \exlevel{Exploration}
  \item Explain the tower property using a population divided into groups. \exlevel{Exploration}
\end{enumerate}
\end{exercisebox}

\chapter{From repetition to regularity}
\chaptertag{Laws of large numbers, the central limit theorem, and sampling}

\begin{question}
Why do averages stabilize, and why do bell curves emerge from non-normal data?
\end{question}

Let $X_1,\ldots,X_n$ be independent copies of a variable with mean $\mu$ and variance $\sigma^2$. Their average is
\[
\bar X_n=\frac1n\sum_{i=1}^nX_i,
\qquad
\E[\bar X_n]=\mu,
\qquad
\Var(\bar X_n)=\frac{\sigma^2}{n}.
\]

\begin{theorem}[Weak law of large numbers]
For every $\varepsilon>0$,
\[
\P(|\bar X_n-\mu|>\varepsilon)\longrightarrow0.
\]
\end{theorem}

With finite variance, Chebyshev gives the proof immediately:
\[
\P(|\bar X_n-\mu|>\varepsilon)
\le \frac{\sigma^2}{n\varepsilon^2}.
\]

\begin{theorem}[Central limit theorem]
If $0<\sigma^2<\infty$, then
\[
\frac{\sqrt n(\bar X_n-\mu)}{\sigma}
\xrightarrow{d}\Normal(0,1).
\]
\end{theorem}

The LLN says where the average goes; the CLT describes the scale and approximate shape of its remaining fluctuations.

\begin{center}
\begin{tikzpicture}[x=.7cm,y=.8cm]
  \foreach \x/\h in {0/1.0,1/3.4,2/1.2,3/3.1,4/.8}{\fill[rose!65] (\x,0) rectangle +(0.7,\h);}
  \node at (1.9,-.6) {one draw};
  \draw[-{Stealth},thick,mint] (4.3,1.8)--(5.5,1.8);
  \foreach \x/\h in {6/.4,7/1.2,8/2.3,9/3.1,10/2.3,11/1.2,12/.4}{\fill[violet!65] (\x,0) rectangle +(0.7,\h);}
  \node at (9.3,-.6) {average of many draws};
\end{tikzpicture}
\end{center}

\section{Sampling distributions}

A \emph{statistic} is a function of the sample. Before data are observed it is random, and its probability law is its sampling distribution. The standard deviation of that distribution is the \emph{standard error}.

For a Bernoulli sample, the sample proportion $\hat p$ has
\[
\E[\hat p]=p,\qquad
\operatorname{SE}(\hat p)=\sqrt{\frac{p(1-p)}n}.
\]

\section{Monte Carlo and bootstrap}

Monte Carlo approximates an expectation by simulating from a known model. The bootstrap approximates a sampling distribution by repeatedly resampling, with replacement, from observed data. The bootstrap is powerful but cannot repair a badly biased or unrepresentative original sample.

\begin{exercisebox}
\begin{enumerate}
  \item Compute the standard error of the mean when $\sigma=12$ and $n=64$. \exlevel{Fluency}
  \item Prove the weak law using Chebyshev's inequality. \exlevel{Reasoning}
  \item Approximate $\P(480\le X\le520)$ for $X\sim\Bin(1000,.5)$ using the CLT and a continuity correction. \exlevel{Fluency}
  \item Explain why quadrupling the sample size halves a standard error. \exlevel{Reasoning}
  \item Simulate sample means from a skewed distribution for $n=1,2,5,30$. Quantify convergence rather than only describing the pictures. \exlevel{Exploration}
  \item Compare Monte Carlo error at $n=100$, $10^4$, and $10^6$. Is the $n^{-1/2}$ rate visible? \exlevel{Exploration}
  \item Give an example where bootstrap resampling cannot reveal a sampling bias. \exlevel{Reasoning}
\end{enumerate}
\end{exercisebox}

\chapter{Learning from a sample}
\chaptertag{Estimation, confidence intervals, and hypothesis tests}

\begin{question}
How can a random sample support a calibrated statement about an unknown population?
\end{question}

An estimator $T=T(X_1,\ldots,X_n)$ targets a parameter $\theta$. Its bias, variance, and mean squared error are
\[
\operatorname{Bias}_\theta(T)=\E_\theta[T]-\theta,
\qquad
\operatorname{MSE}_\theta(T)=\Var_\theta(T)+\operatorname{Bias}_\theta(T)^2.
\]
Unbiasedness is attractive, but MSE captures the bias--variance tradeoff.

\section{Method of moments and likelihood}

The method of moments matches sample moments to model moments. Maximum likelihood chooses the parameter under which the observed data have greatest likelihood.

For Bernoulli observations $x_1,\ldots,x_n$, the likelihood is
\[
L(p)=p^{\sum x_i}(1-p)^{n-\sum x_i},
\]
and the maximizer is $\hat p=\bar x$.

\section{Confidence intervals}

An approximate $95\%$ interval for a mean with known or well-estimated standard error is
\[
\text{estimate}\ \pm\ 1.96\times\text{standard error}.
\]
The confidence level refers to the long-run success rate of the procedure, not a posterior probability that a fixed parameter lies in one realized interval.

\begin{center}
\begin{tikzpicture}[x=1.25cm,y=.55cm]
  \draw[rose,very thick] (4,-.5)--(4,6.5) node[above] {true $\theta$};
  \foreach \y/\a/\b in {0/2.7/4.7,1/3.1/5.0,2/1.8/3.8,3/3.4/5.2,4/2.9/4.4,5/4.2/5.8,6/2.4/4.2}{
    \draw[indigo,thick] (\a,\y)--(\b,\y);
    \fill[indigo] ({(\a+\b)/2},\y) circle (1.6pt);
  }
  \node[left] at (1.7,2) {miss};
\end{tikzpicture}
\end{center}

\section{Tests and $p$-values}

A hypothesis test specifies a null model $H_0$, a test statistic, and a rejection rule. The $p$-value is the probability, assuming $H_0$, of a result at least as incompatible with $H_0$ as the observed one.

It is not the probability that $H_0$ is true. A Type I error rejects a true null; a Type II error fails to reject a false null. Power is the probability of rejecting under a specified alternative.

\begin{warningbox}
Statistical significance is not practical importance. With enough data, a tiny effect can yield a tiny $p$-value; with too little data, an important effect can be missed.
\end{warningbox}

\section{Model checking and design}

Random sampling supports generalization; random assignment supports causal comparison. No inferential formula compensates for selection bias, poorly measured outcomes, hidden multiplicity, or a model contradicted by the data.

\begin{exercisebox}
\begin{enumerate}
  \item Show that the sample mean is unbiased for $\mu$. \exlevel{Reasoning}
  \item Find the Bernoulli maximum-likelihood estimator. \exlevel{Fluency}
  \item Compare two estimators when one has small bias but much smaller variance. \exlevel{Reasoning}
  \item Construct a $95\%$ normal interval for a proportion from 600 observations with 330 successes. \exlevel{Fluency}
  \item Explain the repeated-sampling meaning of that interval without saying there is a $95\%$ chance the fixed parameter is inside. \exlevel{Reasoning}
  \item A test reports $p=.03$. List three claims that do not follow. \exlevel{Reasoning}
  \item Simulate confidence-interval coverage and test power as functions of sample size. \exlevel{Exploration}
  \item Describe how optional stopping or multiple testing can inflate false discoveries. \exlevel{Exploration}
\end{enumerate}
\end{exercisebox}

\chapter{Probability for unknowns}
\chaptertag{An introduction to Bayesian inference}

Chapter~\ref{ch:exchangeability} will revisit this chapter's shared-parameter models from another direction: under infinite exchangeability, a conditional-i.i.d.\ representation is a theorem rather than a starting assumption.

\begin{question}
What changes when an unknown parameter is itself represented by a probability distribution?
\end{question}

Bayesian inference combines a prior distribution $\pi(\theta)$ with a likelihood $p(y\mid\theta)$. Bayes' rule produces the posterior
\[
\pi(\theta\mid y)=
\frac{p(y\mid\theta)\pi(\theta)}{\int p(y\mid u)\pi(u)\,du}
\propto p(y\mid\theta)\pi(\theta).
\]

\section{Beta--binomial updating}

Suppose $p\sim\operatorname{Beta}(\alpha,\beta)$ and, conditional on $p$, $X\sim\Bin(n,p)$. After observing $x$ successes,
\[
p\mid X=x\sim\operatorname{Beta}(\alpha+x,\beta+n-x).
\]
The prior parameters behave like pseudocounts: $\alpha-1$ prior successes and $\beta-1$ prior failures when both exceed one.

The posterior mean is
\[
\E[p\mid x]=\frac{\alpha+x}{\alpha+\beta+n},
\]
a weighted compromise between the prior mean and the sample proportion.

\begin{center}
\begin{tikzpicture}[x=8cm,y=1cm]
  \draw[->,indigo] (0,0)--(1.05,0) node[right] {$p$};
  \draw[amber,thick,domain=.01:.99,samples=70,smooth] plot (\x,{1.3*(1-\x)});
  \draw[mint,very thick,domain=.01:.99,samples=70,smooth] plot (\x,{5.5*\x*\x*(1-\x)});
  \node[amber!80!black] at (.18,1.05) {prior};
  \node[mint!80!black] at (.76,1.05) {posterior};
  \node[below] at (0,0) {0}; \node[below] at (1,0) {1};
\end{tikzpicture}
\end{center}

\section{Credible intervals and prediction}

A $95\%$ posterior credible interval contains $95\%$ of the posterior probability. This is a conditional probability statement about the parameter under the model and prior. A posterior predictive distribution averages future outcomes over posterior uncertainty:
\[
p(y_{\mathrm{new}}\mid y)=\int p(y_{\mathrm{new}}\mid\theta)\pi(\theta\mid y)\,d\theta.
\]

\section{Prior sensitivity and checking}

Priors make assumptions visible; they do not make assumptions disappear. Compare reasonable priors, examine posterior predictive simulations, and ask whether the fitted model reproduces important features of the data.

\begin{bigidea}
Frequentist and Bayesian intervals answer different probability questions. The productive comparison is not ``which formula looks similar?'' but ``what is random, what is conditioned upon, and what operating guarantee or posterior statement is being claimed?''
\end{bigidea}

\begin{exercisebox}
\begin{enumerate}
  \item Update a Beta$(2,2)$ prior after 8 successes in 10 trials. \exlevel{Fluency}
  \item Compute the posterior mean and compare it with the sample proportion. \exlevel{Fluency}
  \item Derive beta--binomial conjugacy by collecting powers of $p$ and $1-p$. \exlevel{Reasoning}
  \item Under a Beta$(\alpha,\beta)$ posterior, find the predictive probability that the next Bernoulli trial succeeds. \exlevel{Reasoning}
  \item Compare Beta$(1,1)$, Beta$(2,2)$, and Beta$(20,20)$ priors for the same small dataset. \exlevel{Exploration}
  \item Explain why a credible interval and confidence interval can have the same numerical endpoints but different interpretations. \exlevel{Reasoning}
  \item Design a posterior predictive check for a coin model that might have time-varying bias. \exlevel{Exploration}
\end{enumerate}
\end{exercisebox}

\section*{Capstone: one question, two inferential languages}

Choose a binomial, count, or measurement dataset. State a generative model, inspect the data, estimate the target, and quantify uncertainty both frequentistically and Bayesianly. End with a model check and a paragraph explaining which conclusions are robust to the modeling choices.

\chapter[Exchangeability]{Exchangeability: when symmetry produces a parameter}
\label{ch:exchangeability}
\chaptertag{A capstone on dependence, learning, and long-run frequency}

\begin{question}
If the order of observations carries no information, must the observations be independent?
\end{question}

No. But for an \emph{infinite} sequence of zeros and ones, a remarkably precise substitute is available: the sequence behaves as though a success probability were chosen once, and independent trials were then run with that probability. This is de Finetti's representation theorem.

This chapter joins three earlier ideas: conditioning can remove dependence, Bayes' rule updates an unknown parameter, and repeated observations reveal a long-run frequency. Read the examples first. Section~\ref{sec:df-proof} is an optional proof, using a birthday-collision bound and one explicitly stated compactness fact from analysis.

\section{Symmetry is weaker than independence}

\begin{definition}
A finite random vector $(X_1,\ldots,X_n)$ is \emph{exchangeable} if permuting its coordinates does not change its joint distribution. An infinite sequence $X_1,X_2,\ldots$ is exchangeable if every finite prefix is exchangeable.
\end{definition}

For a zero--one vector, the definition says that two particular strings with the same number of ones have the same probability. Thus $1100$, $1010$, and $0011$ must be equally likely. It does \emph{not} say that $1100$ and $1111$ are equally likely, or that the coordinates are independent.

Independent and identically distributed (i.i.d.) variables are exchangeable: permuting a product of identical marginal laws changes nothing. Exchangeability also implies identical marginal distributions, because we can swap any two positions. Neither implication reverses in general.

\begin{example}[One coin, chosen once]\label{ex:hidden-coin}
Choose a coin with heads probability $1/4$ or $3/4$, each with probability $1/2$. Keep the chosen coin and toss it repeatedly, independently \emph{given the choice}. Write $X_i=1$ for heads.

For a specified length-$n$ string containing $s$ heads, its probability is
\[
\frac12\left(\frac14\right)^s\left(\frac34\right)^{n-s}
+\frac12\left(\frac34\right)^s\left(\frac14\right)^{n-s}.
\]
Only the count $s$ matters, so the sequence is exchangeable. Nevertheless,
\[
\P(X_1=1)=\frac12,\qquad
\P(X_1=X_2=1)=\frac12\frac1{16}+\frac12\frac9{16}
=\frac5{16}\ne\frac14.
\]
Seeing a head makes the high-bias coin more plausible, and therefore makes another head more likely. The tosses are dependent when the coin's identity is unknown.
\end{example}

\begin{checkpoint}
If we independently choose a fresh coin before \emph{each} toss, the results are instead i.i.d.\ fair bits. Sharing one random choice across all trials is the source of dependence in Example~\ref{ex:hidden-coin}.
\end{checkpoint}

\section{The Bernoulli representation theorem}

\begin{theorem}[de Finetti, Bernoulli case]\label{thm:df}
An infinite zero--one sequence is exchangeable if and only if there is a probability measure $\mu$ on $[0,1]$ such that, for every $n\ge1$ and every $(x_1,\ldots,x_n)\in\{0,1\}^n$ with $s=\sum_i x_i$,
\begin{equation}\label{eq:df}
\P(X_1=x_1,\ldots,X_n=x_n)
=\int_{[0,1]} p^s(1-p)^{n-s}\,\mu(dp).
\end{equation}
The measure $\mu$ is unique. Equivalently, the sequence has the law of conditionally i.i.d.\ Bernoulli$(\Theta)$ trials, where $\Theta$ is chosen once with distribution $\mu$.
\end{theorem}

An exponent zero contributes a factor of one, including at $p=0$ or $1$. The integral averages over a probability \emph{measure}: $\mu$ may have a density, point masses, or both. It need not be a beta distribution. We call $\mu$ the \emph{mixing law} or \emph{de Finetti measure}.

\begin{bigidea}[title={Two levels of randomness}]
First choose $\Theta\sim\mu$. Conditional on $\Theta=p$, generate $X_1,X_2,\ldots$ independently with success probability $p$. Conditional independence concerns the second level; exchangeability describes the joint law after averaging over the first.
\end{bigidea}

The easy direction is worth proving now: every integrand in \eqref{eq:df} depends on the string only through $s$, so every mixture is exchangeable. The difficult direction says that \emph{every infinite exchangeable Bernoulli law} admits such a representation.

Let $S_n=X_1+\cdots+X_n$. There are $\binom ns$ strings with $s$ ones, so
\begin{equation}\label{eq:df-count}
\P(S_n=s)=\binom ns\int_{[0,1]}p^s(1-p)^{n-s}\,\mu(dp).
\end{equation}
The binomial coefficient belongs to the \emph{count} probability, not the probability of one specified ordered string. This distinction is exactly the permutations-and-combinations distinction from Chapter~1.

Ordinary i.i.d.\ Bernoulli$(p_0)$ trials correspond to $\mu=\delta_{p_0}$, a point mass at $p_0$. At the other extreme, choose a fair bit $Z$ and put $X_i=Z$ for every $i$. The sequence is exchangeable, all trials agree, and $\mu=\tfrac12\delta_0+\tfrac12\delta_1$.

\begin{warningbox}[title={What the theorem does not choose}]
Exchangeability specifies a class of models, not a particular prior. Uniqueness means that the \emph{entire infinite joint law} determines $\mu$; it does not mean that finitely many observations determine it. Nor does the representation prove that nature literally selected a hidden coin. Different mechanisms can have exactly the same observable law.
\end{warningbox}

\section{Learning about the shared coin}

Return to Example~\ref{ex:hidden-coin}. After a specified string with $s$ heads in $n$ tosses, Bayes' rule gives
\[
q_{n,s}:=\P\left(\Theta=\frac34\given X_1=x_1,\ldots,X_n=x_n\right)
=\frac{3^s}{3^s+3^{n-s}}.
\]
The probability of heads on the next toss is therefore
\[
\P(X_{n+1}=1\mid X_1=x_1,\ldots,X_n=x_n)
=\frac14(1-q_{n,s})+\frac34 q_{n,s}.
\]
After two heads, $q_{2,2}=9/10$ and the predictive probability is $7/10$. After a head and a tail, $q_{2,1}=1/2$ and it is $1/2$. The order $HT$ versus $TH$ does not affect either answer.

For an arbitrary mixing law, the same calculation gives
\begin{equation}\label{eq:df-predict}
\P(X_{n+1}=1\mid X_1=x_1,\ldots,X_n=x_n)
=\frac{\int p^{s+1}(1-p)^{n-s}\,\mu(dp)}
{\int p^s(1-p)^{n-s}\,\mu(dp)},
\end{equation}
provided the observed string has positive probability. This is the posterior mean of $\Theta$. Conditioning on the count $S_n=s$ gives the same result, since the factor $\binom ns$ cancels.

There is no conflict between exchangeability and learning. Exchangeability says that the \emph{order} of the observed bits carries no additional information beyond their count. It does not say that the count carries no information about future bits.

\section{P\'olya's urn: reinforcement becomes a mixture}
\label{sec:polya}

Start with one red and one blue ball. Draw a ball uniformly, replace it, and add one extra ball of the same color. Repeat. Let $X_i=1$ mean that draw $i$ is red. After $s$ red draws among $n$ draws, the next draw is red with probability
\begin{equation}\label{eq:polya-predict}
\frac{s+1}{n+2}.
\end{equation}
This mechanism seems to remember its past: drawing red makes future reds more likely. Yet its sequence is exchangeable.

\begin{example}[Calculate an ordered path]
The path $RRRB$ has probability
\[
\frac12\cdot\frac23\cdot\frac34\cdot\frac15=\frac1{20}.
\]
The path $RBRR$ has probability $\tfrac12\cdot\tfrac13\cdot\tfrac24\cdot\tfrac35=1/20$ as well. These calculations are not the probability of getting three reds in four draws: there are four possible positions for the blue draw, giving $4/20=1/5$ for that event.
\end{example}

For any specified path with $s$ red and $n-s$ blue draws, the red numerators, in their order of appearance, are $1,2,\ldots,s$; the blue numerators are $1,2,\ldots,n-s$. The denominators are $2,3,\ldots,n+1$. Consequently
\begin{equation}\label{eq:polya-word}
\P(X_1=x_1,\ldots,X_n=x_n)=\frac{s!(n-s)!}{(n+1)!}.
\end{equation}
The order changes where the factors appear, but not their product.

Repeated integration by parts gives the elementary beta integral
\[
\int_0^1p^s(1-p)^{n-s}\,dp=\frac{s!(n-s)!}{(n+1)!}.
\]
Comparing this with \eqref{eq:df} identifies the mixing law: $\Theta$ is uniform on $[0,1]$. Thus two different descriptions generate the same law of observed colors:
\begin{enumerate}
  \item reinforce the color just drawn, using the urn rule;
  \item choose one uniform random bias $\Theta$, then generate conditionally independent Bernoulli$(\Theta)$ draws.
\end{enumerate}
This is equality of joint distributions, not a claim that the physical urn contains a concealed coin.

An especially striking consequence is
\[
\P(S_n=s)=\binom ns\frac{s!(n-s)!}{(n+1)!}=\frac1{n+1},
\qquad s=0,\ldots,n.
\]
The \emph{count} is uniform, even though the individual strings are not. For four draws, all five counts $0,1,2,3,4$ have probability $1/5$. Compare independent fair tosses, whose count probabilities are $1,4,6,4,1$, divided by $16$.

\subsection*{Other starting compositions}
With $a$ red and $b$ blue balls initially, where $a,b$ are positive integers, define the rising factorial $a^{\overline r}=a(a+1)\cdots(a+r-1)$ and $a^{\overline0}=1$. It is different from the falling factorial used to count ordered selections without replacement. The path probability becomes
\[
\frac{a^{\overline s}b^{\overline{n-s}}}{(a+b)^{\overline n}}
=\int_0^1 p^s(1-p)^{n-s}
\frac{p^{a-1}(1-p)^{b-1}}{B(a,b)}\,dp,
\]
where $B(a,b)=\int_0^1p^{a-1}(1-p)^{b-1}\,dp$. The mixing law is Beta$(a,b)$.

The Bayesian model also makes sense for arbitrary real $a,b>0$, whether or not we picture whole balls. After $s$ successes in $n$ trials, the posterior is Beta$(a+s,b+n-s)$ and the predictive success probability is $(a+s)/(a+b+n)$. For the one-red, one-blue urn, this reproduces \eqref{eq:polya-predict}. Chapter~9's beta update now has a concrete reinforcement interpretation.

\Needspace{11\baselineskip}
\section{A law of large numbers with a random limit}
\label{sec:df-limit}

For fixed $p$, the strong law for independent Bernoulli$(p)$ trials says $S_n/n\to p$ with probability one. This means that the set of infinite sample paths where convergence fails has probability zero; it is stronger than convergence in probability in Chapter~7. A short proof of this special strong law appears in Section~\ref{sec:df-proof}.

Apply the strong law conditional on the shared bias. It gives
\begin{equation}\label{eq:df-limit}
\frac{S_n}{n}\longrightarrow\Theta\qquad\text{with probability one}.
\end{equation}
In particular, the mixing law is the distribution of the limiting empirical frequency. The parameter is not merely an algebraic trick: it can be recovered from the infinite sequence itself.

\begin{example}[Same marginal, different limits]
In each of these models $\P(X_i=1)=1/2$:
\begin{center}
\small
\begin{tabularx}{.98\linewidth}{lX}
\toprule
Model & Limit of $S_n/n$ \\
\midrule
Independent fair tosses & $1/2$, with probability one \\
One coin of bias $1/4$ or $3/4$ & $1/4$ or $3/4$, each with probability $1/2$ \\
One-red, one-blue P\'olya urn & A uniform random value in $[0,1]$ \\
One fair bit repeated forever & $0$ or $1$, each with probability $1/2$ \\
\bottomrule
\end{tabularx}
\end{center}
Identical one-trial probabilities tell us very little about the behavior of a long sequence.
\end{example}

We can see the persistent variation without invoking a limit theorem. Conditional independence and the tower property give, for $i\ne j$,
\[
\E X_i=\E\Theta,\qquad
\E[X_iX_j]=\E[\Theta^2],\qquad
\boxed{\Cov(X_i,X_j)=\Var(\Theta)\ge0.}
\]
Also, conditional on $\Theta$, the sample mean has variance $\Theta(1-\Theta)/n$ and mean $\Theta$. Taking expectations in the decomposition into within-group and between-group variation gives
\begin{equation}\label{eq:df-variance}
\boxed{\Var(S_n/n)=\frac{\E[\Theta(1-\Theta)]}{n}+\Var(\Theta).}
\end{equation}
For the two-coin model this is $3/(16n)+1/16$; for the one-red, one-blue urn it is $1/(6n)+1/12$.

Increasing the number of observations reduces the sampling noise \emph{within a run}. It does not remove the variation of the shared bias \emph{between runs}. In the urn example, for instance, $\P(\lim S_n/n>3/4)=1/4$, despite the perfectly symmetric starting urn.

\begin{labbox}[title={Explore: one long run versus many runs}]
Simulate 100 independent copies of the one-red, one-blue urn, each for 5,000 draws. Plot a few running red frequencies; then make a histogram of the 100 final frequencies. Compare with 100 runs of independent fair tosses and with 100 runs of the two-coin model. Predict all three histograms before simulating. A single long trajectory estimates its own shared bias, not the distribution of biases across runs.
\end{labbox}

\Needspace{14\baselineskip}
\section{Why infinity is essential}

\begin{example}[A symmetric pair that cannot continue]
Let $(X_1,X_2)$ be $(1,0)$ or $(0,1)$, each with probability $1/2$. This pair is exchangeable, but
\[
\Cov(X_1,X_2)=0-\frac12\frac12=-\frac14.
\]
It cannot be a mixture of conditionally i.i.d.\ Bernoulli trials, since every such mixture has nonnegative covariance. In fact, it cannot even be the first two coordinates of an exchangeable triple: all three pairs would have to disagree with probability one, which three bits cannot do.
\end{example}

More generally, sample without replacement from $N$ balls, $M$ red and $N-M$ blue. The $N$ color indicators are exchangeable, but for $i\ne j$ and $p=M/N$,
\[
\P(X_i=X_j=1)=\frac{M(M-1)}{N(N-1)},\qquad
\Cov(X_i,X_j)=-\frac{p(1-p)}{N-1}\quad(N>1).
\]
For $0<M<N$, exhausting red balls creates negative dependence. Finite symmetry is real, but it need not be the symmetry of an infinite sequence.

\begin{bigidea}[title={The finite bridge}]
Conditional on its total number of ones, every finite exchangeable zero--one vector is a uniform arrangement of those ones and zeros. It is therefore a mixture of sampling-without-replacement models. For a short sample from a large population, replacement makes little difference. This observation is the bridge to de Finetti's theorem.
\end{bigidea}

\section{Beyond coins, and the limits of symmetry}

The corresponding theorem for an infinite exchangeable sequence of \emph{real-valued} random variables replaces the random number $\Theta$ by a random probability distribution $Q$ on $\R$. Conditional on $Q$, the variables are i.i.d.\ with law $Q$. For Borel sets $A_1,\ldots,A_n$,
\[
\P(X_1\in A_1,\ldots,X_n\in A_n)
=\E\left[\prod_{i=1}^n Q(A_i)\right].
\]
This general result is stated for perspective, not proved here. The real-valued setting matters: unrestricted measurable spaces require additional care.

For example, let $M$ be a random group mean and let $X_i=M+\varepsilon_i$, where the $\varepsilon_i$ are i.i.d.\ $\Normal(0,1)$ and independent of $M$. Conditional on $M$, the observations are i.i.d.\ $\Normal(M,1)$; unconditionally, they are exchangeable and, when $\Var(M)>0$ is finite, have covariance $\Var(M)>0$. Here $Q=\Normal(M,1)$ is a random distribution.

Exchangeability is a modeling assumption to defend. Time trends, different treatment groups, and changing sampling protocols can make permutations inappropriate. Declaring observations exchangeable does not erase such structure. Its value is to make a precise connection: \emph{a symmetry of the observable law can justify a conditional i.i.d.\ representation, without assuming unconditional independence.}

\section{Exercises}
\label{sec:df-exercises}
\begin{exercisebox}[title={Exchangeability practice}]
\begin{enumerate}
  \item In the two-coin model, observe three heads and one tail. Find the posterior probability that the bias is $3/4$ and the predictive probability of another head. Would $HTHH$ give a different answer? \exlevel{Fluency}
  \item Prove that independently selecting a fresh bias from $\{1/4,3/4\}$ before every toss produces i.i.d.\ fair bits. Identify which part of the shared-coin model has changed. \exlevel{Reasoning}
  \item For the one-red, one-blue urn, find the probability of the specified path $RRBRR$, and then the probability of exactly four reds in five draws. Given this path, find the probability of a red next. \exlevel{Fluency}
  \item An urn starts with two red and three blue balls and uses P\'olya reinforcement. After four red and two blue draws, identify the mixing law before observing data, the posterior law, and the next-red probability. \exlevel{Fluency}
  \item Derive \eqref{eq:df-variance}. Under uniform mixing, compute $\Var(S_{10}/10)$ and $\Var(S_{100}/100)$. Explain why neither should be approximated by $1/(4n)$. \exlevel{Reasoning}
  \item Suppose an infinite exchangeable Bernoulli sequence is pairwise independent. Prove that it is i.i.d. What hypothesis prevents this from contradicting Chapter~3's pairwise-independence warning? \exlevel{Reasoning}
  \item Let an exchangeable vector of length $N\ge2$ have common finite variance $v$ and common pairwise covariance $c$. Prove $c\ge-v/(N-1)$. Use this to rule out an exchangeable extension of the always-disagreeing pair to any length $N>2$. \exlevel{Reasoning}
  \item Run the experiment in the Explore box. For each model, state whether final frequencies from many independent runs should cluster near one point, two points, or across an interval. Separate a simulation observation from a theorem that explains it. \exlevel{Exploration}
  \item Use the coupling in Section~\ref{sec:df-proof} to bound the error for \emph{any event} concerning the first five draws from an exchangeable zero--one population of size 1,000. Explain why the same bound need not be informative for 100 draws. \exlevel{Optional proof}
\end{enumerate}
\end{exercisebox}

\subsection*{Sources and a next step}
The representation theorem is due to Bruno de Finetti. For another elementary Bernoulli proof, using moments, see \href{https://arxiv.org/abs/1809.00882}{Werner Kirsch, \emph{An elementary proof of de Finetti's Theorem} (2018)}. Finite exchangeability and sampling approximations are treated systematically in \href{https://doi.org/10.1214/aop/1176994663}{Persi Diaconis and David Freedman, \emph{Finite Exchangeable Sequences} (1980)}.

\clearpage
\section{Optional proof: finite urns and birthday collisions}
\label{sec:df-proof}

Here is a proof of the Bernoulli representation, its uniqueness, and the random-frequency interpretation. The finite approximation is elementary. We isolate the one compactness input rather than conceal it in a limit.

\subsection*{Step 1: condition on the finite total}
Fix $n\le N$ and a particular length-$n$ string with $s$ ones. Conditional on $S_N=k$, exchangeability makes all $\binom Nk$ length-$N$ strings with $k$ ones equally likely, whenever $\P(S_N=k)>0$. Hence
\[
\P(X_1=x_1,\ldots,X_n=x_n\mid S_N=k)
=\frac{\binom{N-n}{k-s}}{\binom Nk}.
\]
A binomial coefficient is zero when its lower argument is outside its allowed range. This is precisely the probability of that ordered color string when sampling without replacement from $k$ ones and $N-k$ zeros.

\subsection*{Step 2: allow replacement and bound the error}
Draw $n$ independent uniform labels $J_1,\ldots,J_n$ from $\{1,\ldots,N\}$. If no label repeats, use these same labels as the without-replacement sample. If a label repeats, generate a new uniform ordered sample of distinct labels for the without-replacement sample. This construction has the correct marginal sampling laws and makes their color strings identical whenever the original labels are distinct.

For each pair $i<j$, $\P(J_i=J_j)=1/N$. The union bound from the birthday problem gives
\begin{equation}\label{eq:df-collision}
\P(\text{a repeated label})\le
\min\left\{1,\frac{\binom n2}{N}\right\}.
\end{equation}
Two coupled random strings that agree except on an event of probability $\varepsilon$ assign probabilities differing by at most $\varepsilon$ to \emph{any} event: their event indicators differ only when the strings differ. With replacement, the colors are independent Bernoulli$(k/N)$. Thus, writing $\mu_N$ for the distribution of $S_N/N$, and averaging over $k$,
\begin{equation}\label{eq:df-approx}
\left|\P(X_1=x_1,\ldots,X_n=x_n)
-\int p^s(1-p)^{n-s}\,\mu_N(dp)\right|
\le\frac{\binom n2}{N}.
\end{equation}
The left-hand probability is fixed as $N$ grows: all prefixes belong to the \emph{same} infinite sequence.

\subsection*{Step 3: take a subsequence of probability laws}
\begin{lemma}[Compactness on a closed bounded interval]\label{lem:df-compact}
Every sequence of probability measures on $[0,1]$ has a subsequence $\nu_j$ and a probability measure $\nu$ on $[0,1]$ such that $\int f\,d\nu_j\to\int f\,d\nu$ for every continuous function $f:[0,1]\to\R$. Also, two probability measures with equal integrals against all such functions are equal.
\end{lemma}

This is the analysis input, also called weak sequential compactness. One route to its proof selects convergent subsequences of distribution functions at each rational point, then uses monotonicity to obtain a limiting distribution function at its continuity points. Compact support prevents probability from escaping to infinity. We use the lemma here without developing that analysis.

Apply the lemma to $\mu_N$. Along one subsequence, $\mu_{N_j}$ converges weakly to some $\mu$. For fixed $n$ and $s$, the function $p^s(1-p)^{n-s}$ is continuous, so \eqref{eq:df-approx} yields \eqref{eq:df}. The same subsequence works for every continuous test function, hence for every finite string. This proves existence.

\subsection*{Step 4: show that the mixing law is unique}
Under \emph{any} candidate mixture, conditional on $\Theta=p$, $S_n\sim\Bin(n,p)$. Chebyshev gives the uniform bound
\[
\P\bigl(|S_n/n-p|>\delta\mid\Theta=p\bigr)
\le\frac{1}{4n\delta^2}\qquad(\delta>0).
\]
For a continuous $f$ on $[0,1]$, let $K=\sup|f|$ and $\omega_f(\delta)=\sup_{|x-y|\le\delta}|f(x)-f(y)|$. Split according to whether $|S_n/n-\Theta|\le\delta$ to get
\[
\left|\E f(S_n/n)-\E f(\Theta)\right|
\le \omega_f(\delta)+\frac{2K}{4n\delta^2}.
\]
First let $n\to\infty$, then $\delta\downarrow0$; uniform continuity makes the right-hand side vanish. Thus the distribution of $S_n/n$ converges weakly to the mixing law. The observable finite-dimensional laws already determine the left-hand side, so two candidate mixing laws have equal continuous-function integrals and are equal by Lemma~\ref{lem:df-compact}.

\subsection*{Step 5: justify the almost-sure random limit}
We first prove the strong law for fixed $p$ without assuming it. For i.i.d.\ Bernoulli$(p)$ trials, Chebyshev along square indices gives
\[
\P\left(\left|\frac{S_{m^2}}{m^2}-p\right|>\varepsilon\right)
\le\frac{1}{4m^2\varepsilon^2}.
\]
The probability that at least one such deviation occurs for some $m\ge M$ is at most $\sum_{m\ge M}(4m^2\varepsilon^2)^{-1}$, which tends to zero. Therefore only finitely many deviations occur with probability one. Intersecting these probability-one conclusions for $\varepsilon=1,1/2,1/3,\ldots$ proves $S_{m^2}/m^2\to p$ almost surely.

Since the partial sums are nondecreasing, for $m^2\le n\le(m+1)^2$,
\[
\frac{S_{m^2}}{(m+1)^2}\le\frac{S_n}{n}
\le\frac{S_{(m+1)^2}}{m^2}.
\]
Both bounds tend to $p$, proving $S_n/n\to p$ almost surely.

Now construct the mixture by choosing $\Theta\sim\mu$ and then the conditionally independent trials. The fixed-$p$ result has conditional probability one for every $p$, so averaging gives $S_n/n\to\Theta$ with probability one. The constructed sequence has the same finite-string probabilities as the original sequence. These determine the probability law on infinite zero--one sequences (including convergence events, which use countable unions and intersections). Thus the original sequence also has an almost-sure empirical limit, distributed as $\mu$.

Taking that limit as the original sequence's $\Theta$ transfers the conditional-i.i.d.\ representation as well: the pair consisting of the sequence and its empirical limit has the same law in the two descriptions. This completes the promised interpretation of the theorem.


\appendix
\chapter{Selected hints and solutions}

\section*{Chapter 1}
For the at-least-one-ace problem, count the complement: $\binom{52}{5}-\binom{48}{5}$. For divisibility by 2, 3, or 5, inclusion--exclusion gives
\[
500+333+200-166-100-66+33=734.
\]
The four-person selection counts from nine students are
$\binom94=126$, $(9)_4=3024$, and $4\binom94=504$.
The repeated-letter count for $\texttt{LEVEL}$ is $5!/(2!2!)=30$.
The coefficient of $x^2$ in $(1+3x)^5$ is $\binom52 3^2=90$.
The length-five strings using all three symbols number
$3^5-3\cdot2^5+3=150$.
For Exercise 3, choose the three letter positions, then fill all positions:
$\binom83 26^3 10^5$.
For Exercises 13--16:
\begin{enumerate}[start=13]
 \item A specified string with three successes has probability
 $(1/3)^3(2/3)^2=4/243$. There are $\binom53=10$ such strings,
 giving total probability $40/243$.
 \item The ordered count is $3^4=81$. The unordered count is
 $\binom{4+3-1}{4}=15$. A selection $AAAA$ has one ordering,
 whereas $AAAB$ has four; the multiplicity is not constant.
 \item There are 11 letters, with four I's, four S's, two P's,
 and one M. The count is $11!/(4!4!2!)=34650$.
 \item Count by subset size to obtain $\sum_k k\binom nk$.
 Alternatively, choose the distinguished member in $n$ ways and
 any subset of the remaining $n-1$ members in $2^{n-1}$ ways.
\end{enumerate}

\section*{Chapter 2}
For the maximum of two dice to equal 4, both dice must be at most 4 but not both at most 3, giving $(4^2-3^2)/36=7/36$. For two uniform points $X,Y$, the band $|X-Y|<.1$ has area $1-2(.9)^2/2=.19$.
For the weighted four-outcome space, the probability of an even outcome
is $(2+4)/10=3/5$, not $2/4$.
On a 30-day calendar, the birthday collision probability is about
$.41356$ for six people and $.53084$ for seven; thus the first size above
$1/2$ is seven. For four people on five days, there are
$5\cdot4\cdot3\cdot2=120$ distinct-birthday lists among $5^4=625$ lists,
so the collision probability is $1-120/625=101/125$.
The quantity $\binom n2/d$ sums probabilities of overlapping events.
It is an upper bound, not necessarily a probability.
For Exercises 12--16:
\begin{enumerate}[start=12]
 \item There are $5\cdot4=20$ equally likely ordered pairs of distinct
 labeled balls. Exactly $3\cdot2=6$ have two reds, giving $3/10$.
 Unordered pairs give $\binom32/\binom52=3/10$ as well.
 \item Choose the fixed pair and derange the remaining three:
 $\binom52D_3/5!=10\cdot2/120=1/6$.
 \item $\P(A_1)=1/4$ and $\P(A_1\cap A_2)=2!/4!=1/12$.
 The product of the single probabilities is $1/16$, so the events
 are dependent.
 \item The preceding recurrence gives
 $(-1)^n=(n-1)D_{n-2}-D_{n-1}$.
 Substitute into $D_n=nD_{n-1}+(-1)^n$ to get the claimed formula.
 Then $D_7=6(265+44)=1854$.
 \item The probability is $\sum_{j=1}^{\infty}4^{-j}
 =(1/4)/(1-1/4)=1/3$.
\end{enumerate}

\section*{Chapter 3}
For the coin problem,
\[
\P(D\mid HHH)=\frac{1\cdot .2}{1\cdot .2+(1/2)^3\cdot .8}=\frac23.
\]
A pairwise-not-mutual example: flip two fair coins; let $A$ be first-heads, $B$ second-heads, and $C$ equal-results.

\section*{Chapter 4}
Exercises 1--18 use discrete variables. Exercises 19--20 belong to
the later continuous section.
\begin{enumerate}
 \item The masses at $0,1,2$ are $1/4,1/2,1/4$. The CDF is $0$
 for $x<0$, $1/4$ for $0\le x<1$, $3/4$ for $1\le x<2$, and
 $1$ for $x\ge2$.
 \item $\E[X]=(1+4+9+16)/10=3$ and
 $\E[X^2]=(1+8+27+64)/10=10$, so $\Var(X)=1$.
 Also $\P(X\ge3)=(3+4)/10=7/10$.
 The equally weighted average $5/2$ uses the wrong probabilities.
 \item Write $m=\E[X]$ and expand:
 $\E[(X-m)^2]=\E[X^2]-2m\E[X]+m^2=\E[X^2]-m^2$.
 \item Let $I_i$ indicate that $i$ is fixed and write $N=\sum I_i$.
 Since $\P(I_i=1)=1/n$, linearity gives $\E[N]=1$ for $n\ge1$.
 No independence is used.
 \item One example pays $999$ with probability $1/1000$ and $-1$
 otherwise. Its mean is zero and its variance is
 $999^2/1000+999/1000=999$. A single play loses one unit with
 probability $999/1000$; a zero mean does not make the payoff certain
 or the chance of a loss small.
 \item Chebyshev gives at most $1/9$. For a fair die the mean is
 $7/2$, the standard deviation is $\sqrt{35/12}$, and every value
 is within $5/2$ of the mean. Since $3\sqrt{35/12}>5/2$, the exact
 probability is zero. The bound is valid but loose.
 \item Equal masses at $-1,1$ have mean zero and variance one.
 Another example has masses $1/4,1/2,1/4$ at
 $-\sqrt2,0,\sqrt2$; its mean is zero and second moment is also one.
 \item The event consists of
 $(4,6),(5,5),(6,4),(5,6),(6,5),(6,6)$, so its probability is $6/36=1/6$.
 \item The masses at $0,1,2$ are $4/9,4/9,1/9$. The CDF is
 $0$ below zero, $4/9$ on $[0,1)$, $8/9$ on $[1,2)$, and $1$
 on $[2,\infty)$. The first two moments are $2/3$ and $8/9$,
 giving variance $8/9-4/9=4/9$. Also $\P(X\ge1)=5/9$.
 \item The values of $Y$ at $X=0,1,2,3$ are $1,0,1,4$.
 Its masses at $0,1,4$ are $3/8,1/2,1/8$, giving mean
 $0(3/8)+1(1/2)+4(1/8)=1$.
 \item The payoff probabilities are $1/2,1/3,1/6$.
 The mean is $-1+1/3+2/3=0$, and the second moment is
 $2+1/3+8/3=5$, so the variance is five.
 The positive-payoff probability is $1/3+1/6=1/2$.
 After a fee of one, the net values are $-3,0,3$: the mean is $-1$,
 variance remains five, and the probability of a positive net payoff
 becomes $1/6$.
 \item At $r=0,1,2,3,4$, the permutation counts are $9,8,6,0,1$,
 respectively, out of $24$. Thus $\E[N]=(8+12+4)/24=1$ and
 $\E[N^2]=(8+24+16)/24=2$, giving variance one.
 Fixing three of four objects forces the fourth to be fixed as well.
 \item There are $\binom{23}{2}=253$ pairs, each matching with
 probability $1/365$, so the mean is $253/365\approx.69315$.
 Multiple pairs can match in one outcome. Markov at threshold one
 gives $\P(C\ge1)\le253/365$, whereas the exact probability is about $.50730$.
 \item The new mean is $2(7/2)-3=4$, variance is
 $4(35/12)=35/3$, and standard deviation is $\sqrt{35/3}$.
 \item Each of $I,J=1-I$ has variance $1/4$, but $I+J=1$
 has variance zero. With an independent second indicator, the
 sum instead has variance $1/4+1/4=1/2$.
 \item The sum has mean seven and variance $35/6$.
 The difference has mean zero and variance $35/6$:
 negation changes the mean's sign but not the variance.
 The average has mean $7/2$ and variance $(35/6)/4=35/24$.
 \item For $n\ge2$, expand
 $N(N-1)=\sum_{i\ne j}I_iI_j$. Each of the $n(n-1)$ ordered
 terms has expectation $1/[n(n-1)]$, so $\E[N(N-1)]=1$.
 Adding $\E[N]=1$ gives $\E[N^2]=2$ and variance one.
 For $n=1$, $N=1$ is constant and its variance is zero.
 \item Write $X-a=(X-m)+(m-a)$, where $m=\E[X]$.
 The cross term has expectation $2(m-a)\E[X-m]=0$.
 Thus $\E[(X-a)^2]=\Var(X)+(m-a)^2$, uniquely minimized by $a=m$.
 \item Normalization gives $c\int_0^2x\,dx=2c=1$, so $c=1/2$.
 The mean is $\int_0^2x^2/2\,dx=4/3$ and the second moment is
 $\int_0^2x^3/2\,dx=2$. The variance is $2-16/9=2/9$.
 \item The density is $3x^2$ on $[0,1]$, zero elsewhere.
 The tail probability is $1-(1/2)^3=7/8$.
 The first two moments are $3/4$ and $3/5$,
 hence the variance is $3/5-9/16=3/80$.
\end{enumerate}

\section*{Chapter 5}
Exercises 1--14 concern discrete families and coupon collection; Exercises 15--17 use
continuous distributions.
\begin{enumerate}
 \item The complement gives
 $1-.7^5-5(.3)(.7)^4=.47178$.
 The mean is $5(.3)=1.5$ and the variance is $5(.3)(.7)=1.05$.
 \item For $0<p<1$ and nonnegative integers $s,t$, divide
 $(1-p)^{s+t}$ by $(1-p)^s$ to get $(1-p)^t$.
 The denominator is positive. If $p=1$, only $s=0$ is admissible;
 conditioning on an event of probability zero is not defined.
 \item Here $np=2$. The exact probability is
 $.998^{1000}\approx.135064522$, compared with
 $e^{-2}\approx.135335283$ for a Poisson variable of mean two.
 The absolute error is about $.000270761$.
 \item The mean is $9/2$, variance $(8^2-1)/12=21/4$,
 and the probability of exceeding $9/2$ is $4/8=1/2$.
 \item The indicator has mean $.2$ and variance $.2(.8)=.16$.
 Ten independent repetitions have a count of mean two and variance $1.6$.
 Recording one indicator ten times gives $10I$, still of mean two but
 variance $100(.16)=16$. Its only possible values are zero and ten.
 \item Without replacement, count each possible number of red balls:
 $\binom22/\binom52=1/10$ for zero reds,
 $\binom31\binom21/\binom52=6/10$ for one, and
 $\binom32/\binom52=3/10$ for two.
 The mean is $6/5$ and variance $9/25$.
 With replacement the masses are $4/25,12/25,9/25$,
 with mean $6/5$ and variance $12/25$.
 \item Examples: a fair die fits discrete uniform, whereas a loaded
 die generally does not. A success indicator fits Bernoulli, whereas
 a count taking values $0,1,2$ does not.
 A count of independent fair tosses fits binomial; repeated copies of
 one toss do not. A uniform subset of a fixed urn fits hypergeometric;
 draws with replacement generally do not.
 Independent repeated trials with a fixed success chance fit geometric;
 changing that chance after each failure generally does not.
 For a Poisson example, take the limiting count of $n$ independent
 opportunities each of probability $\lambda/n$.
 A binomial count with a fixed finite maximum and positive mean cannot
 be exactly Poisson, which has an unbounded set of possible counts.
 \item The roll count is geometric with $p=1/6$, so its mean is six,
 variance is thirty, and $\P(T>5)=(5/6)^5\approx.40188$.
 The failure count $T-1$ has mean five.
 \item $\P(X\ge2)=1-e^{-3}(1+3)=1-4e^{-3}\approx.80085$.
 Also $\E[X(X-1)]=9$, hence $\E[X^2]=9+3=12$ and variance $12-9=3$.
 \item With replacement, the count is $\Bin(4,3/10)$ with
 mean $6/5$ and variance $21/25$.
 Without replacement, the mean is still $6/5$ but the variance is
 $(21/25)(16/19)=336/475$.
 The draws have the same marginal marked probability; without
 replacement, their dependence reduces the variance.
 \item Decompose the collection time into independent geometric stages
 with success probabilities $1,5/6,4/6,3/6,2/6,1/6$.
 The mean is $6H_6=147/10=14.7$ and the variance is
 $36\sum_{j=1}^6j^{-2}-6H_6=3899/100=38.99$.
 When two faces remain, the expected further wait is $6/2+6/1=9$ rolls.
 \item Stop the stage sum after $W_{r-1}$:
 \[
  \E[T_{n,r}]=\sum_{k=0}^{r-1}\frac{n}{n-k}
   =n(H_n-H_{n-r}).
 \]
 The remaining stages have expected total $nH_{n-r}$, even conditional
 on the particular history up to seeing $r$ types. For $r=n$ this is
 zero, using $H_0=0$.
 \item Require $100e^{-m/100}\le.05$, or
 $m\ge100\log2000\approx760.09025$.
 Thus 761 draws suffice. The mean $100H_{100}\approx518.74$ describes
 an average collection time; it does not specify a $95\%$ completion
 deadline. The exponential bound is sufficient and may be conservative.
 \item With a Poisson total, the six face counts are independent
 Poisson variables of mean two. Each is zero with probability $e^{-2}$,
 so the unseen-face count is $\Bin(6,e^{-2})$, with mean $6e^{-2}$.
 The completion probability is $(1-e^{-2})^6\approx.41791$.
 For exactly 12 rolls, inclusion--exclusion instead gives
 \[
  \sum_{j=0}^6(-1)^j\binom6j(1-j/6)^{12}\approx.43782.
 \]
 Those face counts are dependent because they sum to 12. Randomizing
 the total changes the experiment and its completion probability.
 \item The median solves $e^{-\lambda m}=1/2$, giving
 $m=(\log2)/\lambda$.
 \item The standard deviation is three, so standardization gives
 $\P(8<X<13)=\Phi(1)-\Phi(-2/3)$.
 \item For $V$ uniform on $[0,1]$,
 $\E[V]=1/2$ and $\E[V^2]=1/3$, so $\Var(V)=1/12$.
 Write $X=a+(b-a)V$ to obtain mean $(a+b)/2$ and variance $(b-a)^2/12$.
\end{enumerate}

\section*{Chapter 6}
For $X$ uniform on $\{-1,0,1\}$ and $Y=X^2$, symmetry gives $\E X=\E[XY]=0$, hence covariance zero, but knowing $X$ determines $Y$. For the tower property, take the weighted average of group means; each observation receives exactly its population weight.

\section*{Chapter 7}
If $\sigma=12,n=64$, then $\operatorname{SE}(\bar X)=12/8=1.5$. The weak law follows from $\Var(\bar X_n)=\sigma^2/n$ and Chebyshev.

\section*{Chapter 8}
For 330 successes in 600 trials, $\hat p=.55$ and the plug-in standard error is
\[
\sqrt{.55(.45)/600}\approx .0203.
\]
The approximate $95\%$ interval is $.55\pm1.96(.0203)$, or about $(.510,.590)$.

\section*{Chapter 9}
Beta$(2,2)$ updated with 8 successes and 2 failures becomes Beta$(10,4)$, with posterior mean $10/14\approx.714$. The posterior predictive success probability equals the posterior mean $\alpha'/ (\alpha'+\beta')$.

\clearpage
\section*{Chapter 10}
The numbering here matches Section~\ref{sec:df-exercises}.
\begin{enumerate}
  \item The posterior probability of bias $3/4$ is
  $3^3/(3^3+3^1)=9/10$. The predictive chance of heads is
  $(1/4)(1/10)+(3/4)(9/10)=7/10$. Reordering the observed string does
  not change either calculation.
  \item Let the newly chosen biases be independent variables
  $\Theta_1,\ldots,\Theta_n$. For any specified bits,
  \[
  \P(X_1=x_1,\ldots,X_n=x_n)
  =\E\prod_{i=1}^n\Theta_i^{x_i}(1-\Theta_i)^{1-x_i}
  =\prod_{i=1}^n\frac12=2^{-n}.
  \]
  Here the tosses are generated independently given all the biases.
  The expectation factors because the biases are independent; it does
  not factor this way when a single shared $\Theta$ is reused.
  \item The ordered path has probability
  $4!1!/6!=1/30$. There are $\binom54=5$ such paths, so exactly four
  reds has probability $1/6$. After the path, the urn contains five
  red and two blue balls; the next-red probability is $5/7$.
  \item Initially $\Theta\sim\operatorname{Beta}(2,3)$. The posterior
  is Beta$(6,5)$ and the predictive probability is $6/11$.
  The updated urn also contains six red and five blue balls.
  \item Write $Y=S_n/n$ and $m=\E\Theta$. Expand
  $Y-m=(Y-\Theta)+(\Theta-m)$. The cross term has expectation zero
  because $\E[Y-\Theta\mid\Theta]=0$. Therefore
  \[
  \Var(Y)=\E[\Var(Y\mid\Theta)]+\Var(\Theta)
  =\frac{\E[\Theta(1-\Theta)]}{n}+\Var(\Theta).
  \]
  Under uniform mixing, $\E\Theta=1/2$, $\E\Theta^2=1/3$,
  so the two terms are $1/(6n)$ and $1/12$. This gives
  $\Var(S_{10}/10)=1/10$ and $\Var(S_{100}/100)=17/200$.
  The expression $1/(4n)$ wrongly treats the trials as independent
  Bernoulli$(1/2)$ variables; it misses the between-run variance.
  \item Pairwise independence implies $\Cov(X_1,X_2)=0$, hence
  $\Var(\Theta)=0$ by the representation theorem. Thus
  $\Theta=p_0$ almost surely for some constant $p_0$, and
  \eqref{eq:df} is the i.i.d.\ Bernoulli$(p_0)$ product law.
  Infinite exchangeability is the additional hypothesis. An arbitrary
  finite collection of pairwise independent events need not have it.
  \item Nonnegativity of variance gives
  \[
  0\le\Var\left(\sum_{i=1}^N X_i\right)=Nv+N(N-1)c.
  \]
  Divide by $N(N-1)$ to obtain the bound. The always-disagreeing
  pair has $v=1/4,c=-1/4$, so an exchangeable extension would require
  $-1/4\ge-1/[4(N-1)]$, possible only for $N\le2$.
  \item Fair-toss frequencies cluster near $1/2$; shared-coin
  frequencies cluster near $1/4$ and $3/4$; urn frequencies spread
  across $[0,1]$. Finite simulations fluctuate, and 100 runs will not
  give a perfectly flat histogram. The theorem identifies the limiting
  laws as $\delta_{1/2}$, $\tfrac12\delta_{1/4}+\tfrac12\delta_{3/4}$,
  and the uniform law, respectively.
  \item Conditional on the total number of ones, compare sampling
  without replacement with independent sampling at that population's
  success proportion; then average over the total. For any event,
  the absolute probability error is at most
  $\binom52/1000=.01$. With 100 draws the union bound is
  $\binom{100}{2}/1000=4.95$; capped at one, it is only the trivial
  bound. This does not say the approximation error actually equals
  one. The collision argument is useful when $n^2/N$ is small.
\end{enumerate}

\chapter{Reference sheets}

\section*{Counting and overlap corrections}
For $0\le k\le n$, ordered selections without replacement number
$(n)_k=n!/(n-k)!$; unordered selections number
$\binom nk=n!/[k!(n-k)!]$. Ordered $k$-term lists with replacement number
$n^k$. Here $0!=1$.
\begin{align*}
\binom nk&=\binom n{n-k},\\
\binom nk&=\binom{n-1}{k-1}+\binom{n-1}{k},\\
(x+y)^n&=\sum_{k=0}^n\binom nk x^{n-k}y^k,\\
|A\cup B\cup C|&=|A|+|B|+|C|-|A\cap B|-|A\cap C|\\
&\quad-|B\cap C|+|A\cap B\cap C|.
\end{align*}
The same alternating formula holds for event probabilities, without an
independence assumption. Pairwise intersections include the triple region.

\section*{Finite probability spaces and birthdays}
For finite $\Omega$, assign weights $w_\omega\ge0$ summing to one;
then $\P(A)=\sum_{\omega\in A}w_\omega$. Only in the uniform case does
this simplify to $|A|/|\Omega|$.

For $n$ independent uniform birthdays on $d$ days,
\[
\P(\text{some birthday match})=
\begin{cases}
1-\displaystyle\prod_{j=0}^{n-1}(1-j/d),&0\le n\le d,\\
1,&n>d.
\end{cases}
\]
For a designated person among $n\ge1$, the probability of a match with
someone else is $1-(1-1/d)^{n-1}$, taking it to be zero when $n=1$.
For 365 days, the any-pair probability first exceeds $1/2$ at 23 people.

\section*{Derangements and fixed points}
For a uniform permutation of $n$ objects,
\[
 D_n=n!\sum_{k=0}^n\frac{(-1)^k}{k!},\qquad
 \P(\text{exactly }r\text{ fixed points})=\frac{\binom nrD_{n-r}}{n!}.
\]
Here $D_0=1$, $D_1=0$, and $D_n=nD_{n-1}+(-1)^n$.
The error in $D_n/n!\approx e^{-1}$ is less than $1/(n+1)!$;
for $n\ge1$, $D_n$ is the nearest integer to $n!/e$.

\section*{Random variables and threshold events}
The event $\{X\ge a\}$ means $\{\omega:X(\omega)\ge a\}$.
For a PMF, group outcome probabilities with the same value of $X$.
For a density, integrate over the requested range:
\[
 F_X(x)=\P(X\le x),\qquad
 \P(a<X\le b)=F_X(b)-F_X(a),\qquad
 \P(X>a)=1-F_X(a).
\]
For integer-valued $X$ and integer $k$, $\P(X\ge k)=1-F_X(k-1)$.
For a variable with a density, individual values have probability zero.

\section*{Calculus used in the course}
\begin{align*}
 \int_a^b f(x)\,dx&=G(b)-G(a)\quad(G'=f),\\
 e^x&=\sum_{k=0}^{\infty}\frac{x^k}{k!},\\
 \log(1-x)&=-\sum_{k=1}^{\infty}\frac{x^k}{k}\quad(0\le x<1),\\
 \sum_{k=0}^{\infty}r^k&=\frac1{1-r}\quad(|r|<1).
\end{align*}
The birthday estimate is $q_d(n)\approx1-e^{-n(n-1)/(2d)}$;
it is a lower bound when $0\le n\le d$.
See Appendix~\ref{app:calculus} for the derivations, hypotheses,
worked integrations, and quantitative error estimates.

\section*{Core probability identities}
\begin{align*}
\P(A^c)&=1-\P(A),\\
\P(A\cup B)&=\P(A)+\P(B)-\P(A\cap B),\\
\P(A\mid B)&=\frac{\P(A\cap B)}{\P(B)},\\
\P(A)&=\sum_i\P(A\mid B_i)\P(B_i),\\
\P(B_j\mid A)&=\frac{\P(A\mid B_j)\P(B_j)}{\sum_i\P(A\mid B_i)\P(B_i)}.
\end{align*}

\section*{Moments and dependence}
\begin{align*}
\E[aX+bY]&=a\E X+b\E Y,\\
\Var(X)&=\E[X^2]-(\E X)^2,\\
\Cov(X,Y)&=\E[XY]-\E X\E Y,\\
\Var(X+Y)&=\Var X+\Var Y+2\Cov(X,Y),\\
\rho_{X,Y}&=\Cov(X,Y)/(\sigma_X\sigma_Y).
\end{align*}

\section*{Common distributions}
\begin{center}\small
\begin{tabular}{llll}
\toprule
Distribution & Support & Mean & Variance \\
\midrule
Bernoulli$(p)$ & $\{0,1\}$ & $p$ & $p(1-p)$ \\
Binomial$(n,p)$ & $0,\ldots,n$ & $np$ & $np(1-p)$ \\
Geometric$(p)$ & $1,2,\ldots$ & $1/p$ & $(1-p)/p^2$ \\
Poisson$(\lambda)$ & $0,1,\ldots$ & $\lambda$ & $\lambda$ \\
Uniform$(a,b)$ & $[a,b]$ & $(a+b)/2$ & $(b-a)^2/12$ \\
Exponential$(\lambda)$ & $[0,\infty)$ & $1/\lambda$ & $1/\lambda^2$ \\
Normal$(\mu,\sigma^2)$ & $\R$ & $\mu$ & $\sigma^2$ \\
Beta$(\alpha,\beta)$ & $[0,1]$ & $\alpha/(\alpha+\beta)$ & $\alpha\beta/[(\alpha+\beta)^2(\alpha+\beta+1)]$ \\
\bottomrule
\end{tabular}
\end{center}

\section*{Inference checklist}
\begin{enumerate}
  \item What is the target population and parameter?
  \item How were units sampled or assigned?
  \item What model connects observations to the parameter?
  \item What is the estimator and its sampling uncertainty?
  \item Which assumptions are needed for the interval or test?
  \item Is the effect practically meaningful?
  \item What diagnostic or predictive check could reveal model failure?
\end{enumerate}

\clearpage
\section*{Exchangeability and a shared random bias}
For an \emph{infinite} exchangeable zero--one sequence, de Finetti's
theorem gives a unique mixing law $\mu$ on $[0,1]$. Choose
$\Theta\sim\mu$ once; conditional on $\Theta$, the trials are i.i.d.\
Bernoulli$(\Theta)$. For a specified string with $s$ ones,
\begin{align*}
\P(X_1=x_1,\ldots,X_n=x_n)&=\int p^s(1-p)^{n-s}\,\mu(dp),\\
\P(S_n=s)&=\binom ns\int p^s(1-p)^{n-s}\,\mu(dp),\\
\Cov(X_i,X_j)&=\Var(\Theta)\quad(i\ne j),\\
\Var(S_n/n)&=\frac{\E[\Theta(1-\Theta)]}{n}+\Var(\Theta),\\
S_n/n&\longrightarrow\Theta\quad\text{with probability one}.
\end{align*}
The predictive success probability is the posterior mean of $\Theta$.
With a Beta$(a,b)$ mixing law and $s$ successes in $n$ trials,
\[
\Theta\mid\text{data}\sim\operatorname{Beta}(a+s,b+n-s),
\qquad
\P(X_{n+1}=1\mid\text{data})=\frac{a+s}{a+b+n}.
\]
P\'olya's urn, with one-ball reinforcement and $a,b$ initial balls of
the two colors, realizes this model for positive integer $a,b$.
For $a=b=1$, the mixing law is uniform and $S_n$ is uniform on
$\{0,\ldots,n\}$.

\begin{warningbox}
Finite exchangeability alone does not imply a mixture of i.i.d.\
trials. Sampling without replacement can have negative covariance.
Exchangeability also does not choose a particular prior or identify
a unique physical mechanism.
\end{warningbox}

For a prefix of length $n\le N$ from an exchangeable zero--one vector
of length $N$, the finite-mixture approximation using bias $S_N/N$
has absolute probability error at most
$\min\{1,\binom n2/N\}$ for any event on that prefix.
See Chapter~\ref{ch:exchangeability} for the coupling proof.


\backmatter
\chapter*{Further reading}
\addcontentsline{toc}{chapter}{Further reading}
\begin{itemize}
  \item D. Freedman, R. Pisani, and R. Purves, \emph{Statistics}.
  \item J. Blitzstein and J. Hwang, \emph{Introduction to Probability}.
  \item L. Wasserman, \emph{All of Statistics} (for a concise next course).
  \item A. Gelman et al., \emph{Bayesian Data Analysis} (for a more advanced Bayesian treatment).
  \item W. Kirsch, \href{https://arxiv.org/abs/1809.00882}{\emph{An elementary proof of de Finetti's Theorem}}, arXiv:1809.00882 (2018).
  \item P. Diaconis and D. Freedman, \href{https://doi.org/10.1214/aop/1176994663}{\emph{Finite Exchangeable Sequences}}, Annals of Probability \textbf{8}(4) (1980), 745--764.
\end{itemize}

\vfill
\begin{center}\color{muted}\small
Course packet designed as an adaptable teaching draft. Classical results and examples are attributed where appropriate; exposition, exercises, and diagrams are prepared for this packet.
\end{center}

\end{document}
